%  LaTeX support: latex@mdpi.com 
%  For support, please attach all files needed for compiling as well as the log file, and specify your operating system, LaTeX version, and LaTeX editor.

%=================================================================
\documentclass[materials,article,accept,pdftex,moreauthors]{Definitions/mdpi} 

\graphicspath{{figures/}} % path to folder with figures
%\usepackage{caption}
%\usepackage{subcaption}
\usepackage{array}
%\usepackage[flushleft]{threeparttable}
\usepackage{booktabs}
\usepackage{dcolumn}
%\usepackage{longtable}
\usepackage{threeparttablex}
\usepackage{multicol,booktabs}
\setlength{\columnsep}{1cm}
\usepackage{adjustbox}
\usepackage{tabularx}
\usepackage{color}
%=================================================================
% MDPI internal commands
\firstpage{1} 
\makeatletter 
\setcounter{page}{\@firstpage} 
\makeatother
\pubvolume{1}
\issuenum{1}
\articlenumber{0}
\pubyear{2022}
\copyrightyear{2022}
\externaleditor{{Academic Editor: Se-Jin Choi} }
\datereceived{18 October 2022} 
%\daterevised{} % Only for the journal Acoustics
\dateaccepted{1 November 2022} 
\datepublished{} 
%\datecorrected{} % Corrected papers include a "Corrected: XXX" date in the original paper.
%\dateretracted{} % Corrected papers include a "Retracted: XXX" date in the original paper.
\hreflink{https://doi.org/} % If needed use \linebreak
%\doinum{}

%\bibstyle{numeric}
%\usepackage[super]{nth}

%------------------------------------------------------------------
% The following line should be uncommented if the LaTeX file is uploaded to arXiv.org
%\pdfoutput=1

%=================================================================
% Add packages and commands here. The following packages are loaded in our class file: fontenc, inputenc, calc, indentfirst, fancyhdr, graphicx, epstopdf, lastpage, ifthen, lineno, float, amsmath, setspace, enumitem, mathpazo, booktabs, titlesec, etoolbox, tabto, xcolor, soul, multirow, microtype, tikz, totcount, changepage, attrib, upgreek, cleveref, amsthm, hyphenat, natbib, hyperref, footmisc, url, geometry, newfloat, caption

%=================================================================
%% Please use the following mathematics environments: Theorem, Lemma, Corollary, Proposition, Characterization, Property, Problem, Example, ExamplesandDefinitions, Hypothesis, Remark, Definition, Notation, Assumption
%% For proofs, please use the proof environment (the amsthm package is loaded by the MDPI class).

%=================================================================
% Full title of the paper (Capitalized)
\Title{Dynamics 
%MDPI: Dear Authors,
%(1) Please proofread the article format according to the comments left and pay attention to the highlighted parts.
%(2) Please do not delete the comments during your proofreading, since we need to check point by point after receiving your proofed version.  - Answer: agree.
%(3) Please just modify or confirm or give feedback with track change Function. - Answer: agree.
%(4) English editing and layout have been complete, please do not revert the changes. - Answer: agree.
%(5) In order to process this paper smoothly, please check if anything is missing in the main text and offer responses to these comments. Thank you very much for your cooperation. - Answer: agree.
%MDPI: This is an important note to let you know that to protect the privacy of the author's contact information, we will only display the corresponding authors' contact information on the published paper on and after 4 August 2022. However, we hope the contact information of the other authors can also be confirmed or corrected during the proofreading stage, which will be recorded in our database to furnish the author's integrated publication history and also contribute to our future communications. - Answer: agree.
% --------------------------------- ANSWER  --------------------------------- %
%-- Answer:  Ok, we have read and understand all your notes. Our replies are marked as "-- Answer: agree (or reply)... " just next to your comments. After we agree with your comments, the yellow highlighting is deleted, but your comments remainedm as you requested.
of Strength Gain in Sandy Soil Stabilised with Mixed Binders Evaluated by Elastic P-Waves during Compressive Loading}

% MDPI internal command: Title for citation in the left column
\TitleCitation{Dynamics of Strength Gain in Sandy Soil Stabilised with Mixed Binders Evaluated by Elastic P-Waves during Compressive Loading}

% Author Orchid ID: enter ID or remove command
\newcommand{\orcidauthorA}{0000-0002-0577-9936} % Add \orcidA{} behind the author's name
\newcommand{\orcidauthorB}{0000-0002-5759-1089} % Add \orcidA{} behind the author's name

% Authors, for the paper (add full first names)
\Author{Per %MDPI: Please carefully check the accuracy of names and affiliations. Please note that authorship can not be changed after manuscript has been accepted. - Answer: agree, name "Per" is correct.
 Lindh $^{1,2}$\orcidA{} and Polina Lemenkova $^{3,}$*\orcidB{}}

%\longauthorlist{yes}

% MDPI internal command: Authors, for metadata in PDF
\AuthorNames{Per Lindh, Polina Lemenkova}

% MDPI internal command: Authors, for citation in the left column
\AuthorCitation{Lindh%MDPI: Please carefully check the accuracy of names. -- Answer: Yes, names are correct.
, P.; Lemenkova, P.}
% If this is a Chicago style journal: Lastname, Firstname, Firstname Lastname, and Firstname Lastname.

% Affiliations / Addresses (Add [1] after \address if there is only one affiliation.)
\address{%
$^{1}$ \quad Swedish Transport Administration, P.O. Box 366, SE-201 23, Neptunigatan 52, Malm\"{o}, Sweden. per.lindh@trafikverket.se 
%MDPI: Please add post code. -- Answer: checked, added exact address with post code.
%MDPI: This email address is different from the one submitted online at susy.mdpi.com  	(per.a.lindh@gmail.com). Please confirm which one is correct. -- Answer: the official emails are correct: we have added 2 official email of the 1st author for each of his affiliations: 1) per.lindh@trafikverket.se for affiliation-1 and 2) per.lindh@byggtek.lth.se for affiliation-2. Both are correct, please include his both official emails.
 \\
$^{2}$ \quad Division of Building Materials, Department of Building and Environmental Technology, Lunds Tekniska H\"{o}gskola (LTH), Faculty of Engineering, Lund University, %MDPI: Please check that the address information is complete. The provided information should be arranged from subordinate to superior. Same question for affiliation 2. -- Answer: checked, correct.
 P.O. Box 118 %MDPI: Please also add this information for affiliation 1 or we will remove this information. -- Answer: corrected (arranged from subordinate to superior). -- Answer: information added in affiliation 1, correct.
, SE-221-00 Lund, Sweden. per.lindh@byggtek.lth.se \\
$^{3}$ \quad Laboratory of Image Synthesis and Analysis, Building L, \'{E}cole Polytechnique de Bruxelles, Brussels Faculty of Engineering, Campus de Solbosch, Universit\'{e} Libre de Bruxelles, ULB--LISA CP165/57, Avenue Franklin D. Roosevelt 50, B-1050 Brussels, Belgium} %-- Answer: correct.

% Contact information of the corresponding author
\corres{Correspondence: polina.lemenkova@ulb.be; Tel.: +{32-471860459}} 

% Abstract (Do not insert blank lines, i.e. \\) 
\abstract{This paper addresses the problem of stabilisation of poor subgrade soil for improving its engineering properties and stiffness. The study aim is to evaluate the effects from single and mixed binders on the gain of strength in sandy soil over the period of curing. We propose an effective non-destructive approach of using P-waves for identifying soil strength upon stabilisation. The growth of strength and stiffness is strongly dependent on time of curing and type of the stabilising agents which can include both single binders and their blended mixtures. The diverse effects from mixed binders on the properties of soil were evaluated, compared and analysed. We performed {the} experimental trials of five different binders for stabilisation of sandy soil using cement, lime, Ground Granulated Blast Furnace Slag (GGBFS), energy fly ash and bio fly ash. The methodology included soil stabilisation by binders during a total period of 90 days, strength test for the Unconfined Compressive Strength (UCS) and seismic tests on the stabilised samples. The dynamics of soil behaviour stabilised by different binders for days 7, 14, 28 and 90 was statistically analysed and compared. The optimisation of binder blending has been performed using mixture simplex lattice design with three binders in each case as independent variables. Using P-waves naturally exploited strength characteristics of soil samples and allowed us to compare the effects from the individual and blended binders over the complete period of curing with dominating mixes. The results indicate that strength growth in stabilised soil samples is nonlinear in both time and content of binders with dominating effects from slag which contributed the most to the compressive strength development, followed by cement. }

% Keywords
\keyword{civil engineering; stabilisation; compressive strength; soil; cement; slag; fly ash; lime; seismic waves} 

% The fields PACS, MSC, and JEL may be left empty or commented out if not applicable
\PACS{81.40.Cd; 81.40.Ef; 62.20.Qp; 83.50.Xa; 45.70.Mg; 92.40.Lg; 81.40.Lm; 62.20.M--%MDPI: Please check if it is correct here. %-- Answer: correct. PACS classification: "62.20.M- Structural failure of materials" with a long 'dash' after 'M'.
}
\MSC{76Axx; 74Exx; 74Fxx}
\JEL{Q00; Q01; Q24; Q55; Q56}

%%%%%%%%%%%%%%%%%%%%%%%%%%%%%%%%%%%%%%%%%%
\begin{document}

\section{Introduction}

Soil stabilisation which can be defined as an improvement and refinement of soil parameters through modification of its engineering properties is currently being used in many of the state-of-the-art high-level civil engineering tasks such as construction works, maintenance of built environment, including structural components of buildings, roads, bridges, airports, etc. The~goal of laboratory experiments on soil stabilisation, given a context of road and highway construction, is to determine the compaction characteristics and the strength of the stabilised material. Stabilisation of soil with traditional binders is used in a wide variety of engineering applications with the major aim at improving soil properties prior to construction works in road engineering \citep{Park,Glossop,Kallenetal2016,Abdullah,Kallenetal2014,Sudhakar,Nordmark}. Cement and lime are the oldest, fundamental and most well-known traditional binders, continuously used in civil engineering \citep{Mitchell,Chen,Hilt,Kaniraj}. Therefore, their applications for soil stabilisation have been researched and reported for decades as a straightforward, versatile and traditional method of soil improvement \citep{Markwick,Abbey,Aldaood,Bell,Dahlinetal1999}. 

Many types of binders have been tested in industrial works and proposed in the literature that can be broadly classified into two groups---traditional binders and novel mixes of blended binders. Although~traditional binders, such as cement or lime, remain the most commonly used additives in soil stabilisation, the~selection of binders becomes more complicated owing to the recent issues. These include negative environmental impacts from cement, the~increased amount of the industrial by-products which should be utilised and the fabrication of novel stabilising agents which have improved geotechnical properties in terms of engineering performance and better suited to the ecological challenges. For~instance, such cases arise when geotechnical works are planned on the expansive or clayey soils that have specific properties~\cite{Nelson}. Moreover, compared to the traditional binders, blended binders have been shown to produce more consistent effects on soil stabilisation by adaptively balancing between the effective refinement of soil parameters and economic costs of works. Blended binders, combined from several materials of various additives can respond to such needs and requirements, since they are better suited to different types of soil and regional environmental conditions, thus considering {real project~conditions.

In this regard, blended binders can be considered as an advanced method of soil stabilisation which has been shown to perform better than traditional binders in the state-of-the-art techniques of soil stabilisation \citep{Jegandanetal2010,Jafer2017}. Although~it derives soil stabilisation results by improving its parameters, which leads to an increased strength of foundation and durability of structures, the~use of traditional binders is restricted in the environmental and technical performance. The~experimental use of the alternative and blended binders allows for application of more effective approaches, both in economical and engineering \mbox{aspects~\cite{Mavroulidou},} which is valuable, due to the high cost of geotechnical works on \mbox{road construction. }

Various attempts have been made to understand how the addition of of blended binders can enhance soil properties and increase its strength, and~in which proportions binders should be taken \citep{AlSwaidani,Fasihnikoutalab,Kogbara,Verhasselt,Sherwood,Chaddock}. Indeed, mixed binders have the potential to improve soil stabilisation through more appropriate workflow of the experiments aimed at reducing soil permeability, compressibility, deformation and settlement. For~instance, the~use of blended binders supports the effectiveness of clayey soil hardening over time of curing~\cite{Wangetal2015,Gandhi,LindhLemenkova2021a}. 

Specifically for Sweden, the~performance of single and mixed binders was investigated using deep mixed method \citep{Ahnberg1995,HolmAhnberg}. The~most well-known application of blended binders in Sweden is the depth stabilisation technique, which follows the long-term tradition of soil treatment in the country~\cite{Bjerrum}. Others include using blended binders for shallow soil stabilisation. Stabilisation has also gained increased use in infrastructure projects for treatment of the contaminated soil \citep{Sorengard,LindhLemenkova2022b}. Further, blended binders perform well compared to single binders, because~such mixtures are better suitable on soils which with varied grain size \citep{ElMahdi,Syed}. Since the response of such soils varies to the different types of binders, blended mixes can be applied for testing soils with diverse grain content, ranging from a fine-grained clay to a coarse gravel. Moreover, the~use of mixed binders in Sweden can help make cemented soil structure estimates consistent with the  {in situ}   conditions of real projects, such as road engineering in cold environment with cases of ground~freezing. 

As the areas of applications increase along with the use of complex mixtures of binders, so do the requirements for the optimisation of soil. In~the post-treatment of the contaminated soil, it is common for chemically active additives to be used to increase the effect of stabilisation with regard to the reduction of the environmental risk. It is also common for cement or slag to be replaced by other binders, such as fly ash or other residual products, with~aim to reduce costs of engineering works \citep{Arulrajah,LindhLemenkova2022a,Jaferetal2018}. Moreover, in~response to the increased societal demands for environmentally sustainably solutions in engineering works, e.g.,~reduced CO$_2$ emissions, the~demands on the environmental impact from binders during the process of soil stabilisation are also increasing. Overall, these requirements and new challenges in road construction industry lead to the need to gain access to novel methods and materials that enable optimisation of soil treatment that would take into account both technical performance, economics and environmental aspects which are to be considered in the final~project.

The Uniaxial Compressive Strength (alternatively: Unconfined Compressive Strength) (UCS) test is a straightforward and relatively easy method to evaluate the enhancement of strength and stiffness in stabilised soil over time. Because~it well reflects both compaction properties and gained strength, it is one of the most frequently used and commonly acceptable methods to evaluate the engineering properties of the stabilised soil, as~continuously reported in earlier papers~\cite{Eskisar,Feng,Rodriguez,Christensen}. The~examples of using the UCS tests are diverse and include, for~instance, the~improved methods testing bearing capacity and settlement of cohesive soil in the foundation of buildings~\cite{Alshkane}. Some more advanced methods include predictive modelling by programming techniques to estimate the UCS~\cite{Onyelowe}, using artificial intelligence methods \citep{Tabarsa,Sihag,Liuetal2015} and machine learning applications \citep{Furtney,Sihag} to UCS estimation. Almost all previous UCS-based soil testing techniques have been restricted to existing workflow using traditional binders. However, the~major disadvantage of the UCS method consists in its destructive nature, since it crushes the tested sample during the~experiment.

Therefore, the~proposed framework uses seismic measurements which alleviates the above difficulties and disadvantages of of the UCS by measuring soil strength using elastic P-waves through non-destructive measurements. The~seismic testing presents an alternative method of evaluation of stiffness and strength parameters of soil stabilised by various percentages of binders which can be performed using measured velocities of elastic pressure waves (P-waves). The~core principle of the evaluation of velocities of P-waves consists in the physical theory of waves and elastic properties of soils as a porous media. Thus, measuring P-wave speed enables to estimate the level of stiffness and strength in a soil sample, as~used in various geotechnical application \citep{Gheibi,Jiangetal2021,Jiangetal2017,LindhLemenkova2021b,Lacidogna}. The~technical approach of this method consists in the propagation of elastic waves that penetrate the specimens of soil materials tested at different curing periods. A~quantitative assessment of the P-wave velocity gives the value of the soil strength and stiffness due to the correlation between these parameters, which is reported in relevant case studies \citep{LindhLemenkova2022c,LinWang2020,Kurtulus,Ryden}. Seismic methods applied to measuring the properties of soil demonstrated superiority over common methods due to its non-destructive nature, which enables to measure samples as many times as needed, but~requires complex tools and advanced methods of data processing \citep{Sarro,Shietal2017}. 

This paper deals with experimental trials of five different binders for stabilisation of clayey soil: cement, lime, GGBFS and two types of fly ash (energy and bio fly ash). The~combinations of blended binders were determined using special techniques of simplex design, aimed at optimising the proportions of the stabilising agents. The~methodology comprises soil treatment by blended binders during the overall period of 90 days, and~geotechnical tests according to the standards of the Swedish Geotechnical Institute (SGI), Link\"{o}ping: strength test for the UCS and seismic measurements (P-wave speed) of the stabilised specimens. The~dynamics of soil behaviour stabilised by different binders was statistically analysed and compared on the 7th, 14th, 28th and 90th days, respectively. The~development in P-wave velocity was compared both for single and blended binders over the complete period of curing, i.e.,~90 days. The~results are summarised in tables and visualised on plots in relevant chapters below. Our main contribution can be summarised as follows: (1) the approach of P-wave tests that performs non-destructive measurements of soil strength in selected samples during the period of curing; (2) stabilisation using both single and blended binders; (3) using statistical simplex design approach as an associated combinatorial structure and factor experiment for optimisation of binder quantity and water ratio; and (4) the use of statistical methods for data analysis based on measurements of P-wave velocities as a function of curing time for individual and blended binary~binders. 

%%%%%%%%%%%%%%%%%%%%%%%%%%%%%%%%%%%%%%%%%%
\section{Materials and~Methods}

\subsection{Materials}


\subsubsection{Specimens}
Soil used in this study was excavated from the selected places in southern Sweden. After~collecting and refining, the~specimens were processed in the laboratory of SGI. Afterwards, the~soil grain size was examined according to the grain size distribution. Based on the examined soil, its structure is dominated by the two major types: sand and gravel, according to the grain size (Figure \ref{fig:01}). 

\begin{figure}[H]
         		\includegraphics[width=10.0cm]{Fig_01.jpg}
     \caption{{Grain} size distribution for soil used in this~study.} \label{fig:01}
\end{figure}

The soil used in this study has a coarse structure and includes fine to  medium silty sand and gravel inclusions obtained in the study area of southern Sweden. According to the Unified Soil Classification System (USCS) the sandy soil samples can be classified as follows: clayey sands, sand-clay mixtures (SC) up to 10\% of samples, silty sands, sand--silt mixtures (SM) and poorly graded sands (SP) (from 10\% to 18\% of samples), and~well-graded sands, gravelly sands (SW) (18--25\% of sample) and belongs to group sands with half or more of coarse fraction, Figure~\ref{fig:01}. The~first group of gravel specimens (25--50\%) is presented by clayey gravels and gravel-sand-clays (GC) and silty gravels, gravel-sand-silt mixtures (GM), followed by medium gravel (over 50\%), characterised by poorly-graded gravels, gravel-sand mixtures (GP), see Figure~\ref{fig:01}. 

The grain size of sand is ranging between the 0.06 to 20 mm (sand) and for 2.0 to \mbox{31.5 mm} for the gravel. The~ballast material (support layer) was weighed in three batches of 1700 g each. Water was measured in measuring glasses, so that the corresponding 102 g per batch weighed had a water ratio of 6\% before the addition of binder and 5.8\% when the binder was included. To~ensure the significance level of the results, the~recipes of binder-soil mixtures were based on the previous experience of similar tests in SGI in combination with the statistical experimental tests (simplex). The~statistical trial design minimises the number of trials that need to be performed with the desired level of significance. Thus, for~instance, for~the five different binders, we fabricated approximately 100 soil~specimens. 

The specimens were stabilised by binders and stored for 90 days of curing period to achieve final strength, even for the slow-hardening components. During~this time, the~velocity of the P-waves and in selected cases also the speed of the shear wave have been measured in the material at different times. The~long curing period of up to 90 days, was selected, because~it is the one of the key factors that affect soil stabilisation. After~the 90th day of storage, selected samples were reduced in size twice, after~which the plastic tube was removed. Long period of curing ensures the improvement of strength which tends to increase with time due to the bonded particles in a soil-binder mixture. Moreover, the~development of soil-lime pozzolanic reaction increases over~time. 

\subsubsection{Binders}
The stabilisation of soil by binders was aimed to improve its geotechnical characteristics prior to construction works. Specifically these include the following goals: to reduce and prevent settlement of soil as a road basement; to increase soil strength; to improve bearing capacity of soil; to lower soil permeability; to prevent shrinkage when undergoing changes in moisture (water) content; to control volume changes in soil caused by the thaw-freeze effects, which is a subject to environmental and climate effects, as~often the case for Sweden. To~this end, we used five different binders for soil stabilisation: two traditional binders (ordinary Portland cement (OPC) and lime) and three alternative binders (slag GGBFS, bio fly ash and energy fly ash), originated from the SCA Lilla Edet and the from coal combustion, certified by International Organization for Standardization (ISO).

The GGBFS slag is of type Merit 5000, which is an additive material for cement and concrete made from dried Hyttsand. Merit 5000 has been used in combination with OPC due to similar effects on strength development in soil} \cite{SSAB}. The~bio fly ash originated as a by-product from the CHP power plants that use only biomass as fuel where pure bio fly ash is separated from the gas in electrofilters before the gas is directed to the chimney. The~bio fly ash was used as an alternative binder and a component to binder blends in this study. The~first fly ash origin is from Svenska Cellulosa Aktiebolaget (SCA) Lilla Edet and the second is from coal combustion (ISO certified). The~bearing layer gravel (0--18) is used as a ballast material. The~choice of these binders was approved for the three~reasons. 

\begin{enumerate}
  \item First, these binders were selected based on the grain size of the specimens, to~be suitable to the material used in a real project, where a maximum particle size reached 20 mm (from fine sand to course gravel). 
  \item Second, it was necessary to use a material with small variation in physiochemical properties, which would minimise the fluctuations in final stabilisation results. 
  \item Third, the~additives should be inert and not contribute to the strength gain in combination with any other binder. In~the first step, sample specimens were fabricated and stabilised using different binder recipes. 
\end{enumerate}

Detailed schedule of the binder recipes for each soil specimen is presented in Table \ref{tabA1}. 

\subsection{Workflow}

The workflow included the evaluation of the effects from single and mixed binders on the gain of strength in soil specimens over time of curing. The~experimental trials included treatment of soil by five different binders used for stabilisation of clayey soil: OPC, lime, GGBFS, energy fly ash and bio fly ash. The~optimisation included simplex design combinatorics method for correct selection of binder combinations and amounts. The~experimental planning aimed at reducing the financial risks of the project, since large amounts of soil had to be stabilised in a real project (several hundreds of tons of soil). 

The optimisation of binder blending was performed in the laboratory of SGI using modelling by mixture simplex lattice design with three binders in each case as independent variables. The~performed methodology include soil treatment by 5 types of binders during 90 days, and~geotechnical tests according to the SGI standards: strength test for the Unconfined Compressive Strength (UCS) by the SIS standard SS-EN 13286-41 and seismic tests on the stabilised samples using ICP Accelerometer for measuring P-wave velocity. The~dynamics of soil behaviour stabilised by different binders for days 7, 14, 28 and 90 was statistically analysed and compared. The~curves of strength gain in soil specimens have been compared using various ratios of binders (double and triple combinations of stabilising agents). Furthermore, the~variation in water ratio in specimens before and after the complete curing period were statistically~assessed.

\subsubsection{Uniaxial Compressive Strength (UCS)}
The UCS tests were carried out on stabilised soil using the 10-ton pressure press, Figure~\ref{fig:02}a. The~tests were run under the deformation control following the existing standard methodology SS-EN 13286-41~\cite{SIS2003,SIS2018}. The~UCS was performed to achieve the maximum axial compressive stress that a specimen can bear under zero confining stress. In~such a way, the~specimens were tested using the available equipment at SGI (Figure \ref{fig:02}a,b for schematic and real view, respectively). A~workflow methodology is based on the the Swedish Institute for Standards (SIS) \cite{SIS2003,SIS2018}. The~breaking load reached 60 s, which was achieved at a deformation rate of 3--4 mm/s, Figure~\ref{fig:02}b. No visible connections between the bearing boundaries and fracture surfaces were noted on any of the tested specimens. The~specimens exhibited the hourglass-shaped fracture surfaces that were completely independent of the layer~boundaries. 

\begin{figure}[H]
     \centering
     \begin{tabular}{cc}
		\includegraphics[width=6.3cm]{Fig_02a.jpg} &
	\includegraphics[width=6.0cm]{Fig_02b.jpg} \\	
	(\textbf{a}) Scheme of the equipment & (\textbf{b}) Real case use of the UCS test device\\
	\end{tabular}
        \caption{{Compaction}  equipment for soil stabilisation (\textbf{a}); Specimen that reached failure (\textbf{b}).}
        \label{fig:02}
\end{figure}



\subsubsection{P-Wave~Measurements}
During the curing period of 90 days, measurements of the resonant frequency of the elastic waves were performed on soil specimens on days 7th, 14th, 28th and 90th using free-free resonate column approach. The~velocity of the pressure waves (P-waves) was measured in specimens on control days during curing period, to~evaluate the dynamics of strength gain. The~P-wave velocities of the stabilised soil were received using the Integrated Circuit Piezoelectric (ICP) Accelerometer (Model Nr. 352B10), a~lightweight ceramic shear response device designed by the PicoCoulomB (PCB) {Piezotronics}  \url{https://www.pcb.com}. The~workflow included the following steps. An~impulse force was excited in a tested specimen, and~the velocity of pressure wave was recorded by the device. The~first arrival time of wave was estimated to compute the travel time of the compression wave. The~velocity of the pressure waves was computed using travel time of the pulse propagated through the soil specimens. The~data were then transferred to the computer for \mbox{statistical processing.
}


\subsubsection{Determination of Water~Content}

The determination of water content has been performed using the existing methodology~\cite{SIS2014}. The~variation in water ratio between the specimens before and after the curing time of 90 days is shown in Figure~\ref{fig:03} which illustrates this difference. Here the blue colour shows water content before curing, while the red colour shows water content after the 90~days of curing. Since samples were paraffined, no water remained in soil samples, but~the difference in water ratio can be explained by chemical reactions during~curing. 

\begin{figure}[H]
         		\includegraphics[width=11.0cm]{Fig_03.jpg}
     \caption{Histogram of water content in soil~specimens.} \label{fig:03}
\end{figure}

Target water ratio for the stabilised material was 5.8\% when binder was included. The~histogram in Figure~\ref{fig:03} illustrates the mean pre-storage of 5.89\% with a standard deviation of 0.307; after curing the mean was 5.37 and the standard deviation $-$0.359, which indicates that the target value was achieved. The~variation in water content was not greater than expected, which depends, i.a. %MDPI: should it be `i.e.'? -- Answer: no, this is correct, it means Latin abbreviation "inter alia", i.a. - meaning 'among other things', 'among others'.
, on~the composition of the material used for water ratio determination. In~samples with a low water ratio, there have been coarser materials and vice~versa for those with a high water ratio. This is because coarser grains cannot bind as much water as several smaller grains together can. The~measurement was performed as a quality control to discard any statistical outliers. For~these reasons, two samples were discarded due to the incorrect values of water quotas. The~samples reported in Figure~\ref{fig:03} only include those accepted after statistical~evaluation.
 
\subsubsection{Freeze-Thaw~Tests}
The freeze-thaw tests were performed according to the Swedish technical standard SS-EN 137244~\cite{SIS2005}. In~this method, sample specimens were frozen for a total period of 56 days. The~amount of material was measured on days {7}th, {14}th, {28}th, {42}th and {56}th. The~temperature varied and demonstrated the dependance on curing time. No extra water was added during these cycles to ensure the objectivity of the experiment. The~samples used for the freeze-thaw tests were selected from the original soil mass in connection with the UCS tests. Therefore, the~specimens that did not hold together could not be tested, for~instance, the~specimens stabilised by energy fly ash or bio flay ash in large amounts which did not stick together. These tests showed that all  tested samples had a high frost~resistance. 

The processes of freezing and thawing, as~common environmental and climate induced phenomena in Sweden, often lead to the increase in volume of soil and reduction of strength. Hence, it is important to consider durability of soil as the final goal with regard to the improved soil properties prior to the road construction. At~the same time, freeze-thaw test results may fluctuate and differ, which depends on the actual field conditions. This is owing to the placement of the stabilised soil in the in situ conditions below the road pavement. As~a result, changes in moisture condition in soil sample may vary from saturated to dry. Likewise, stabilised soil samples may not be frozen when saturated due to the local climate specifics in southern~Sweden. 

\subsection{Concepts}
The conceptual optimisation of the workflow using experimental design is an important part of geotechnical works, as~it enables to minimise the number of tests. In~view of the large amounts of material that should be tested (hundreds of tons), the~optimisation of work is necessary to plan the workload in a more economically effective way. To~assess the impacts from the individual and mixed binders, a~simplex mixture design was chosen following the existing methodology~\cite{MyersMontgomery}, Figure~\ref{fig:04}. 

\begin{figure}[H]
\hspace{-10pt}
\begin{tabular}{cc}
       	  	\includegraphics[width=5.2cm]{Fig_04a.jpg} &
       	  	\includegraphics[width=7.5cm]{Fig_04b.jpg} \\
       	  	(\textbf{a}) {Two} different binders & (\textbf{b}) Three different binders.\\
	\end{tabular}
        \caption{Application %MDPI: Please change the hyphen (-) into minus sign ($-$, "U+2212"). e.g., "-1" should be "$-$1". -- Answer: ok we agree, but there is no hyphen in this phrase "Application of simplex lattice design showing possible binder combinations". We don't understand what is needed.
 of simplex lattice design showing possible binder~combinations.}
        \label{fig:04}
\end{figure}

The role of the trial planning consists of a systematic methodology aimed at finding the optimal and minimal necessary content of the tested binder products which should be added into the soil mixture. The~traditional method to optimising binder mixtures uses the approach which consists in one factor tested at a time. This strategy enables only one parameter to be changed, while keeping all the other parameters as constant values. However, by~this approach, there is a risk of missing both positive and negative interactions between the constituent components of the soil-binder mixture. Neither does it allow a simultaneous optimisation which could be based on the analysis of both technical and environmental requirements of the stabilised~soil.

The decision regarding optimal binder mixture is based on a series of different random trials where various factors can be varied in an unsystematic way. The~mentioned above problems can be eliminated if the statistical trial planning is used to support the optimisation of binder mixtures. Thus, trial design ensures that there is a systematic evaluation of binder mixtures, which is achieved through simultaneous considering of the influence of various factors on the final results~\cite{Montgomery}. Specifically for soil testing, one can measure the effects of various binder combination on the evolution of soil strength over the period of curing. By~working systematically, the~possibility of making a good optimisation of blended mixtures of binders is increased with regard to both technical performance, economic gain and environmental~impact.  

Laboratory tests on binder proportions and combinations often entail higher costs due to many trials in the initial stage. Using this method, high total cost of the geotechnical works on soil stabilisation intended for road construction can be significantly reduced. Since both the choice of binder components and binder quantities largely affect project economy, the~price of works can be reduced in relation to the real needs. Therefore, to~optimise binder mixtures we used modelling trial design techniques. The~approach can be divided into a mixing part and a process control part. Mixture optimisation included the mutual proportions of different binders in the final mixture, that is, the~ratios of blended binders~\cite{Lindh2001}. The~process control optimisation stage focuses on how much binder should be used based on the properties of the base soil material, e.g.,~how much water should be added to balance the ratio and, if~necessary, what is the level of contamination as environmental process~parameters.

\subsubsection{Optimization of Binder~Mixture}
For the two binders, the~optimisation process is a relatively simple and straightforward process. In~this case, all the possible binder combinations can be placed on a straight line that can be described according to Equation~(\ref{eq:1}), Figure~\ref{fig:04}a. As~long as $x < 1$ or $y < 1$, the~mixture consists of two parts. If~$x=1$ or $y=1$, the~mixture consists of a single \mbox{binder component. }
\begin{equation}\label{eq:1}
y~=~1 - x
\end{equation}

Possible binder combinations with two different binders are represented by a straight line where the endpoints signify one binder, 1.0~=~100\%, i.e.,~a one-component mixture. Respectively, adding the {2}nd binder automatically reduces the amount of the {1}st binder, so that together they always constitute a total amount of 100\%, see Figure~\ref{fig:04}a. However, for~mixtures consisting of three or more binders, the~situation becomes more complicated, as~the optimal final mixture is hidden among the infinite number of possible mixture combinations. The~experimental planning handles this case using the simplex method, initially developed in statistical works \citep{Montgomery,MyersMontgomery}. The~general aim of simplex method is to evaluate different types of mixtures with three or more constituent components, which is a very common optimisation scenario in real case situations in industry. Figure~\ref{fig:04}b illustrates a graph representing all possible mixing ratios for the three components A, B and C, based on the simplex method. The~corners represent 100\% of a single binder. The~border lines represent mixtures between the two different binders, according to Equations~(\ref{eq:2})--(\ref{eq:4}).
\begin{equation}\label{eq:2}
y~=~1 - x	
\end{equation}
\begin{equation}\label{eq:3}
z~=~1 - y
\end{equation}
\begin{equation}\label{eq:4}
x~=~1 - z
\end{equation}

The plane bounded by the edge lines consists in the infinitely many combinations of all three binders which can be experimentally mixed and tested. All the possible binder combinations with the three different binders (A, B and C) are represented by a plane bounded by the three straight lines (Figure \ref{fig:04}b). Here each vertex in the plane represents a single binder, i.e.,~1.0~=~100\% of a single component~\cite{MyersMontgomery}. Depending on the scale and amount of works chosen for simplex method, both linear and interaction effects can be identified. These include, for~instance, the~interaction effect which occurs when two binder components together produce a significantly higher or lower effect on soil strength than their mutual proportions indicated earlier. The~level of details, however, places different demands on the total number of tested trials. Therefore, the~amount of the experimental trails should always be reasonably high. The~optimal combination of binders is required for the effective soil stabilisation in geotechnical and economic~aspects.

\subsubsection{Process Control~Optimization}
The process control optimisation examines how much binder must be added based on the properties of the base material, e.g.,~water ratio, organic carbon content or impurity content, the~process parameters of soil. In~the experimental planning, this optimisation has been carried out using the factorial experiment. Figure~\ref{fig:05}a illustrates a simple $2^2$-factor experiment where two factors are tested at the two levels of various binders $2^2$. This method is used to evaluate how two different factors affect the final result. The~factor test can be performed with a higher resolution, e.g.,~by increasing the number of levels to three when the $2^3$-factor is obtained. Similar to the simplex method, higher resolution means that the number of trial experiments increases. The~choice of the resolution basically depends on the real project situation (amount, type and quality of soil that should be stabilised) and which interactions are expected to occur in soil treatment~\cite{Lindh2004}.




\subsubsection{Total Binder~Optimization}
In a total optimisation stage, simplex experimental trials from the mixture optimisation are combined in a factorial trials with process variables. In~this case, the~combined experimental set-up then becomes advanced, as~illustrated schematically in Figure~\ref{fig:05}b. The~advantage of this approach is that we find an optimum value both in terms of binder combination and binder quantity, based on the important process parameters. Otherwise, if~we only investigate the optimum of stabilising mixture, based on only binder combination, we may miss the unforeseen and unwanted influence from the external process parameters. The~reverse applies if we only optimise binder amount based on the process~parameters. 

\begin{figure}[H]
\hspace{-10pt}
\begin{tabular}{cc}
		\includegraphics[width=6.2cm]{Fig_05a.jpg} &
	\includegraphics[width=7.0cm]{Fig_05b.jpg}
	\\	
	(\textbf{a}) & (\textbf{b})
	\end{tabular}
        \caption{\hl{An example} %MDPI: Please replace with a sharper image for subfigure b.
        of a $2^2$-factor experiment with binder quantity and water ratio as factors (\textbf{a}); An example of a $2^2$-factor trial in combination with simplex trials (\textbf{b}).}
        \label{fig:05}
\end{figure}

This type of the total optimisation gives us a robust basis for being able to vary both the binder combinations and the amount of binder in a project, without~risking either final quality or project economics. As~a result, the~combined methodology adopted \mbox{from~\cite{Smith}} provides both optimal binder mix and optimal binder quantity in relation to the \mbox{process parameters.}

%%%%%%%%%%%%%%%%%%%%%%%%%%%%%%%%%%%%%%%%%%
\section{Results and~Discussion}

The dynamics of the P-wave velocity, corresponding to soil strength as a function of curing time is shown in Figure~\ref{fig:06}. For~the interpretability reasons, only the values of the soil specimens stabilised by one type of binder are shown. For~soils that are sandy, added energy fly ash and bio fly ash at all ages of curing the sandy soil exhibited almost no effects on the compressive strength (brown lines in Figure~\ref{fig:06}) with very low values of P-waves (<250 m/s).  Moreover, the~dynamics remain stable for this case and do not change significantly for the period of 90 days. In~contrast, stabilisation of soil by bio fly ash shows a clearly visible sharp increase in the P-wave velocity starting from day 28 from 350 m/s up to 3000 and 3100 for the two different types of bio flay ash (blue line in Figure~\ref{fig:06}). 


Strong effects from added OPC  on strength gain at an early age on the specimen are notable. The~reaction of the two types of cement with different soil specimens varied, but~the general trend is comparative: the values of P-waves start from 2400 and 265 m/s for the two types of cement and gradually continue to gain in strength until the P-waves are recorded at speed of 2600 and 3000 m/s, respectively. With~sandy soil, the~reaction of Portland cement type Cem II/A-V with the sandy fraction is comparable for both types. The~gap in the P-wave speed for lime (purple line in Figure~\ref{fig:06}) and bio fly ash (green line in Figure~\ref{fig:06}) between 28 and 90 days is caused by the disturbances during cutting and demolding process of the~specimens. 



\begin{figure}[H]
         	\includegraphics[width=13.0cm]{Fig_06.jpg}
     \caption{{P-wave} velocity as a function of curing time for individual~binders.} \label{fig:06}
\end{figure}





This disturbance has not affected tested specimens with high strength (cement and slag), i.e.,~the bond between the particles of the specimens stabilised with lime, energy ash and bio fly ash was so small that they broke during handling. Furthermore, lime, energy ash and bio fly ash have actively bounded the aggregates and for the small elongation levels that arise in the connection with seismic measurements, the~bonds remain intact because of the larger elongation levels that arise during cutting and demolding, so the bonds are damaged. The~strength of soil stabilised by lime reached its peak on day 28th when values of the P-waves were recorded as 950 and 1250 m/s for the two sets of measurements. Afterwards, the~gain in strength decreased and values of P-waves dropped until 500 m/s. 

The velocity of the P-waves propagating specimens stabilised by binary binders are shown in Figure~\ref{fig:07}. Specifically, it shows samples stabilised by cement (as a reference), slag and blended binary mixtures in combinations for binders as slag/lime and lime/slag with proportions: 0.67/0.33 for both combinations. A~comparison between the slag and binary binders containing lime and slag shows a faster curing process when slag is combined with lime. Figure~\ref{fig:07} furthermore shows that the two binary binders lime/slag and slag/lime have equally high or higher P-wave velocity on day 90 as cement. This can be compared with the values of the compressive strength reported for the day {90}th for the strength measurements of soil stabilised by cement, lime and GGBFS (Tables 7 and 8). Thus, from~Table 7 we can see that in terms of strength, cement and slag/lime in proportions 0.67/0.33 are equivalent, because~soil stabilised by binary binders lime/slag (0.67/0.33) has a clearly lower strength. This phenomenon may have been owing to the nonlinearities of the stabilised soil material, i.e.,~elongation-dependent~stiffness.


The results from the measured P-wave velocities indicating strength gain (UCS) are reported in Tables 1--8 for days 7, 14, 28 and 90, as~well as the remaining values for the same days, respectively. The~analysis of Table~\ref{tab:1}, which shows measured values of P-waves on the day 7th of curing period of stabilised soil, gives the following remarks. The~presence of cement and energy fly ash (combination AD) influences the stabilisation characteristics of a soil with P-waves (coefficient 2408.8) by T~=~47 $^\circ$C. Likewise, the~addition of lime and slag (combination BC) into soil results in the increase of P-wave coefficient up to 5474.4. The~compound of lime and slag acts as an accelerator for soil hardening and in their presence increase the strength~gain. 


\begin{figure}[H]
         		\includegraphics[width=13cm]{Fig_07.jpg}
     \caption{{P-wave} velocity as a function of curing time for individual and blended binary~binders.} \label{fig:07}
\end{figure}

Correspondingly, adding slag (GGBFS) and bio fly ash (combination CE) leads to the increase of bonding soil particles, which is also important in providing higher particle friction and better packing. As~a result, soil strength increases which is reflected in the higher speed of the elastic waves (P-wave coefficient~=~2453.1). The~measurements were done by T~=~47 $^\circ$C for all the data in Table~\ref{tab:1}. Depending on the amount of added slag (type Merit 5000) and Portland cement, bonding soil particles generally increases the UCS values compared to the one in the fly ashes, which demonstrated poor or low effects. The~improvement is introduced by glueing of soil particles with cementitious binder at the points of contact by stabilising~agents. 

\begin{table}[H]
\renewcommand{\arraystretch}{1.3}
\caption{The %MDPI: Please add an explanation for color in the table footer. If the color is unnecessary, please remove it. Same question for all the following tables. -- Answer: agree and explained: "The~color-marked rows indicate a significance level of values below the \emph{p}~=~0.05" in the table footer. the same for all the following tables. The color is necessary.
 effects from various binders on soil stabilisation on day 7th of curing time, measured by P-wave velocities.}\label{tab:1}
\newcolumntype{C}{>{\centering\arraybackslash}X}
\begin{tabularx}{\textwidth}{CCCCCCC}
\toprule
			\textbf{Factors} & \textbf{P-Wave Coeff. }& \textbf{Std. Err.} & \textbf{T~=~47} \boldmath{$^\circ$}\textbf{C} & \textbf{\emph{{p}}} 
			 & \boldmath{$-$}\textbf{95.\%} \textbf{Cnf. Limt} & \textbf{+95.\%}\textbf{ Cnf. Limt} \\ \noalign{\hrule height .5pt} 
       \rowcolor{green!10} (A) Cement & 2538.3 & 217.223 & 11.68544 & 0.000000 & 2101.4 & 2975.35 \\ \noalign{\hrule height .5pt} 
       \rowcolor{green!10} (B) Lime & 563.5 & 217.223 & 2.59394 & 0.012613 & 126.5 & 1000.46 \\ \noalign{\hrule height .5pt} 
        (C) GGBFS & 277.5 & 217.223 & 1.27759 & 0.207668 & $-$159.5 & 714.52 \\ \midrule
        (D) E\_Fly ash & 174.4 & 217.223 & 0.80280 & 0.426134 & $-$262.6 & 611.38 \\ \noalign{\hrule height .5pt} 
      \rowcolor{green!10}  (E) B\_Fly ash & 1192.9 & 217.223 & 5.49141 & 0.000002 & 755.9 & 1629.86 \\ \noalign{\hrule height .5pt} 
        AB & 730.2 & 972.388 & 0.75095 & 0.456429 & $-$1226.0 & 2686.40 \\ \midrule
        AC & 1345.3 & 972.388 & 1.38346 & 0.173061 & $-$610.9 & 3301.45 \\ \noalign{\hrule height .5pt} 
      \rowcolor{green!10}  AD & 2408.8 & 972.388 & 2.47718 & 0.016896 & 452.6 & 4364.97 \\ \noalign{\hrule height .5pt} 
        AE & $-$1599.1 & 972.388 & $-$1.64453 & 0.106740 & $-$3555.3 & 357.07 \\ \noalign{\hrule height .5pt} 
    \rowcolor{green!10}    BC & 5474.4 & 972.388 & 5.62986 & 0.000001 & 3518.2 & 7430.60 \\ \noalign{\hrule height .5pt} 
        BD & 108.9 & 972.388 & 0.11203 & 0.911278 & $-$1847.3 & 2065.13 \\ \midrule
        BE & $-$962.7 & 972.388 & $-$0.99003 & 0.327230 & $-$2918.9 & 993.50 \\ \midrule
        CD & $-$18.4 & 972.388 & $-$0.01892 & 0.984982 & $-$1974.6 & 1937.79 \\ \noalign{\hrule height .5pt} 
    \rowcolor{green!10}    CE & 2453.1 & 972.388 & 2.52278 & 0.015086 & 496.9 & 4409.31 \\ \noalign{\hrule height .5pt} 
      DE & $-$201.2 & 972.388 & $-$0.20692 & 0.836967 & $-$2157.4 & 1754.99 \\ 
    
        
      
      
      \bottomrule
     \end{tabularx}
        \end{table}

        
        \begin{table}[H]
        \renewcommand{\arraystretch}{1.3}
        \ContinuedFloat
\caption{{\em~Cont.} \label{tab:1}}
\newcolumntype{C}{>{\centering\arraybackslash}X}
\begin{tabularx}{\textwidth}{CCCCCCC}
\toprule
			\textbf{Factors} & \textbf{P-Wave Coeff. }& \textbf{Std. Err.} & \textbf{T~=~47} \boldmath{$^\circ$}\textbf{C} & \textbf{\emph{{p}}} & \boldmath{$-$}\textbf{95.\%} \textbf{Cnf. Limt} & \textbf{+95.\%}\textbf{ Cnf. Limt} \\ \noalign{\hrule height .5pt} 
			   ABC & $-$7209.2 & 6848.162 & $-$1.05272 & 0.297852 & $-$20,985.9 & 6567.54 \\ \midrule
        ABD & 4591.6 & 6848.162 & 0.67049 & 0.505827 & $-$9185.1 & 18,368.35 \\ \midrule
        ABE & 5768.0 & 6848.162 & 0.84227 & 0.403902 & $-$8008.7 & 19,544.73 \\  \midrule			
			ACD & 13,776.3 & 6848.162 & 2.01168 & 0.050007 & $-$0.4 & 27,553.01 \\ \midrule
        ACE & 5997.5 & 6848.162 & 0.87578 & 0.385601 & $-$7779.2 & 19,774.24 \\ \noalign{\hrule height .5pt} 
      \rowcolor{green!10}  ADE & $-$23,505.9 & 6848.162 & $-$3.43244 & 0.001257 & $-$37,282.6 & $-$9729.20 \\
        \noalign{\hrule height .5pt} 
        BCD & $-$3246.1 & 6848.162 & $-$0.47402 & 0.637683 & $-$17,022.9 & 10,530.59 \\ \midrule
        BCE & $-$13,506.6 & 6848.162 & $-$1.97230 & 0.054476 & $-$27,283.3 & 270.12 \\ \midrule
        BDE & 5566.1 & 6848.162 & 0.81279 & 0.420437 & $-$8210.6 & 19,342.84 \\ \midrule
        CDE & 407.6 & 6848.162 & 0.05952 & 0.952793 & $-$13,369.1 & 14,184.31 \\ \midrule
        AB(A-B) & 1814.7 & 2064.817 & 0.87886 & 0.383951 & $-$2339.2 & 5968.55 \\ \midrule
        AC(A-C) & 1048.7 & 2064.817 & 0.50789 & 0.613907 & $-$3105.2 & 5202.57 \\ \midrule
        AD(A-D) & $-$3526.9 & 2064.817 & $-$1.70809 & 0.094218 & $-$7680.8 & 626.99 \\ \noalign{\hrule height .5pt} 
       \rowcolor{green!10}    AE(A-E) & 5648.6 & 2064.817 & 2.73563 & 0.008756 & 1494.7 & 9802.45 \\ \noalign{\hrule height .5pt} 
        BC(B-C) & $-$2271.6 & 2064.817 & $-$1.10012 & 0.276881 & $-$6425.4 & 1882.32 \\ \midrule
        BD(B-D) & $-$265.0 & 2064.817 & $-$0.12835 & 0.898421 & $-$4418.9 & 3888.86 \\ \midrule
        BE(B-E) & 3986.1 & 2064.817 & 1.93050 & 0.059591 & $-$167.7 & 8140.01 \\ \midrule
        CD(-D) & $-$277.3 & 2064.817 & $-$0.13429 & 0.893750 & $-$4431.2 & 3876.60 \\ \midrule
        CE(C-E) & 75.9 & 2064.817 & 0.03677 & 0.970822 & $-$4077.9 & 4229.81 \\ \midrule
        DE(D-E) & 528.2 & 2064.817 & 0.25583 & 0.799196 & $-$3625.6 & 4682.13 \\ \bottomrule
\end{tabularx}
\footnotesize{Factors %MDPI: We removed “Notes for table 1”. Please confirm. Same question for all the following tables. -- Answer: agree, correct (the same for all the following tables).
 are the combinations of binders. All the factors are included in the cubic model. The~color-marked rows indicate a significance level of values below the \emph{p}~=~0.05. The~coefficients for recoded computations are as follows: Var.: Vp\_rod\_7; R-sqr~=~0%MDPI: Please confirm added 0. -- Answer: agree, correct.
.8806; Adj: 0.7942 (Resultat\_rev091123.sta). Five-factor mixture design; Mixture total~=~1, %MDPI: dot removed. Please confirm. same question for all the following table footers. -- Answer: agree, correct (the same for all the following tables).
 82 runs. DV: Vp\_rod\_7; MS Residual~=~95,430.29. Std. Err. is a standard error (standard deviation of sample distribution).}
\end{table}
 
Similarly, the~combination AE (cement and bio flay ash) also increased the soil strength, as~the stabiliser content was increased in specimens. This results in the P-wave coefficient reached 5648.6 m/s. As~for triple blends, upon~exposure to added slag GGBFS and cement into binder blends, the~compressive strength of a stabilised soil is gradually increased after the day 28th. This proves the important role of cement and lime for stabilisation. Thus, the~rate of increase may vary for blends ACD---cement/slag/fly ash (P-wave coefficient~=~13,776.3 m/s), ABE--cement/lime/bio fly ash (P-wave coefficient~=~5768.0), combination ABD---cement/lime/energy fly ash (P-wave coefficient~=~4591.6). Mechanically, a~sandy soil system, when stabilised with either Portland cement, GGBFS or lime, provides improved shear and compressive strength, expressed in the UCS values, as~it is also reflected in the P-waves velocity. Table~\ref{tab:2} shows the remaining notable effects on day 7th of curing period after backward elimination. Here the combination AE (cement/bio fly ash) demonstrated the highest values of P-wave coefficient (5213.9) followed by the blend BC (lime/slag) with the P-wave coefficient of~4311.8.

\begin{table}[H]
\caption{{Remaining} significant effects on day {7}th after the backward~elimination. \label{tab:2}}
\setlength{\cellWidtha}{\columnwidth/7-2\tabcolsep+0.2in}
\setlength{\cellWidthb}{\columnwidth/7-2\tabcolsep+0.0in}
\setlength{\cellWidthc}{\columnwidth/7-2\tabcolsep-0.1in}
\setlength{\cellWidthd}{\columnwidth/7-2\tabcolsep-0.1in}
\setlength{\cellWidthe}{\columnwidth/7-2\tabcolsep-0.0in}
\setlength{\cellWidthf}{\columnwidth/7-2\tabcolsep-0.0in}
\setlength{\cellWidthg}{\columnwidth/7-2\tabcolsep-0.0in}
\scalebox{1}[1]{\begin{tabularx}{\columnwidth}{>{\PreserveBackslash\centering}m{\cellWidtha}
>{\PreserveBackslash\centering}m{\cellWidthb}
>{\PreserveBackslash\centering}m{\cellWidthc}
>{\PreserveBackslash\centering}m{\cellWidthd}
>{\PreserveBackslash\centering}m{\cellWidthe}
>{\PreserveBackslash\centering}m{\cellWidthf}
>{\PreserveBackslash\centering}m{\cellWidthg}}
\toprule
        \textbf{Factors} & \textbf{P-Wave Coeff.} & \textbf{Std. Err.} & \textbf{t (72)} & \emph{\textbf{{p}}} & \boldmath{$-$}\textbf{95.\%} \textbf{Cnf. Limt} & \textbf{+95.\%} \textbf{Cnf. Limt} \\ \midrule
        (A) Cement & 2643.3 & 138.311 & 19.11103 & 0.000000 & 2367.6 & 2919.0 \\ \midrule
        (B) Lime & 601.7 & 133.783 & 4.49725 & 0.000026 & 335.0 & 868.4 \\ \midrule
        (C) GGBFS & 497.6 & 152.147 & 3.27048 & 0.001649 & 194.3 & 800.9 \\ \midrule
        (D) E\_Fly ash & 272.1 & 134.969 & 2.01612 & 0.047519 & 3.1 & 541.2 \\ \midrule
        (E) B\_Fly ash & 893.4 & 139.921 & 6.38482 & 0.000000 & 614.4 & 1172.3 \\ \midrule
        AD & 2789.3 & 809.587 & 3.44535 & 0.000955 & 1175.4 & 4403.2 \\ \midrule
        BC & 4311.8 & 788.740 & 5.46665 & 0.000001 & 2739.4 & 5884.1 \\ \midrule
        CE & 2368.8 & 790.947 & 2.99492 & 0.003762 & 792.1 & 3945.5 \\ \midrule
        ADE & $-$27,704.7 & 6364.530 & $-$4.35298 & 0.000044 & $-$40,392.1 & $-$15,017.2 \\ \midrule
        AE(A-E) & 5213.9 & 2104.584 & 2.47741 & 0.015582 & 1018.5 & 9409.3 \\ \bottomrule
        \end{tabularx}}
    \footnotesize{A~significance level below the \emph{p}~=~0.05 is noted for all the recorded values in Table~\ref{tab:2}. Annotation for Table~\ref{tab:2}: Coeffs (recoded comps); Var.: Vp\_rod\_7; R-sqr~=~0.8014; Adj: 0.7765 (Resultat\_rev091123.sta). Five-factor mixture design; Mixture total~=~1, 82 Runs DV: Vp\_rod\_7; MS Residual~=~103,610.4. Factors stand for combination of binders.}
\end{table}




The analysis of Table~\ref{tab:3} demonstrates measured values of P-waves on the day 14th of curing period. Here we can note that the addition of single binders contributed to the increase of P-wave coefficient, which is attributed to the development of its strength. Specifically, added Portland cement resulted in the P-waves coefficient raising up to 2645.6, followed by bio fly ash (P-wave coefficient~=~1476.8) and lime (P-wave coefficient~=~781.8), while the effects from the energy fly ash and GGBFS are similar (211.8 and 208.5, respectively). In~contrast, blended binders have a more notable effects on the formation of bonded particles, which strengthen a soil that is stabilised with the cement/slag GGBFS (blend AC, P-wave~=~5229.2), lime/GGBFS (blend BC, P-wave~=~7920.3) and slag/energy fly ash (blend CE, P-waves~=~4819.4). As~a result, this produces a higher strength, which is mirrored in the measured speed of the elastic~P-waves. 
 
\begin{table}[H]
 \renewcommand{\arraystretch}{1.3}
\caption{{The effects} from various binders on soil stabilisation on day {14}th of curing time, measured by P-wave~velocities.}\label{tab:3}
\setlength{\cellWidtha}{\columnwidth/7-2\tabcolsep+0.2in}
\setlength{\cellWidthb}{\columnwidth/7-2\tabcolsep+0.0in}
\setlength{\cellWidthc}{\columnwidth/7-2\tabcolsep-0.1in}
\setlength{\cellWidthd}{\columnwidth/7-2\tabcolsep-0.1in}
\setlength{\cellWidthe}{\columnwidth/7-2\tabcolsep-0.0in}
\setlength{\cellWidthf}{\columnwidth/7-2\tabcolsep-0.0in}
\setlength{\cellWidthg}{\columnwidth/7-2\tabcolsep-0.0in}
\scalebox{1}[1]{\begin{tabularx}{\columnwidth}{>{\PreserveBackslash\centering}m{\cellWidtha}
>{\PreserveBackslash\centering}m{\cellWidthb}
>{\PreserveBackslash\centering}m{\cellWidthc}
>{\PreserveBackslash\centering}m{\cellWidthd}
>{\PreserveBackslash\centering}m{\cellWidthe}
>{\PreserveBackslash\centering}m{\cellWidthf}
>{\PreserveBackslash\centering}m{\cellWidthg}}
\toprule
	\textbf{Factors} & \textbf{P-Wave Coeff. }& \textbf{Std. Err.} &\textbf{ t(47)} & \textbf{\emph{p}} & \boldmath{$-$}\textbf{95.\% Cnf. Limt }& \textbf{+95.\% Cnf. Limt} \\ \noalign{\hrule height .5pt}
         \rowcolor{green!10} (A) Cement & 2645.6 & 202.397 & 13.07116 & 0.000000 & 2238.4 & 3052.7 \\ \noalign{\hrule height .5pt}
       \rowcolor{green!10}  (B) Lime & 781.8 & 202.397 & 3.86294 & 0.000342 & 374.7 & 1189.0 \\ \noalign{\hrule height .5pt}
        (C) GGBFS & 208.5 & 202.397 & 1.03027 & 0.308156 & $-$198.6 & 615.7 \\ \midrule
        (D) E\_Fly ash & 211.8 & 202.397 & 1.04639 & 0.300731 & $-$195.4 & 619.0 \\ \noalign{\hrule height .5pt}
        \rowcolor{green!10} (E) B\_Fly ash & 1476.8 & 202.397 & 7.29644 & 0.000000 & 1069.6 & 1883.9 \\ \noalign{\hrule height .5pt}
        AB & 828.6 & 906.019 & 0.91455 & 0.365095 & $-$994.1 & 2651.3 \\ \noalign{\hrule height .5pt}
       \rowcolor{green!10}  AC & 5229.2 & 906.019 & 5.77161 & 0.000001 & 3406.5 & 7051.9 \\ \noalign{\hrule height .5pt}
       \rowcolor{green!10}  AD & 2644.7 & 906.019 & 2.91908 & 0.005374 & 822.1 & 4467.4 \\ \noalign{\hrule height .5pt}
        AE & $-$1669.3 & 906.019 & $-$1.84248 & 0.071717 & $-$3492.0 & 153.4 \\ \noalign{\hrule height .5pt}
       \rowcolor{green!10}  BC & 7920.3 & 906.019 & 8.74183 & 0.000000 & 6097.6 & 9742.9 \\ \noalign{\hrule height .5pt}
        BD & 351.7 & 906.019 & 0.38813 & 0.699673 & $-$1471.0 & 2174.3 \\ \midrule
        BE & $-$980.5 & 906.019 & $-$1.08217 & 0.284699 & $-$2803.1 & 842.2 \\ \bottomrule
 \end{tabularx}}
  \end{table}




\begin{table}[H]\ContinuedFloat
 \renewcommand{\arraystretch}{1.3}
\caption{\emph{Cont.}}\label{tab:3}
\setlength{\cellWidtha}{\columnwidth/7-2\tabcolsep+0.2in}
\setlength{\cellWidthb}{\columnwidth/7-2\tabcolsep+0.0in}
\setlength{\cellWidthc}{\columnwidth/7-2\tabcolsep-0.1in}
\setlength{\cellWidthd}{\columnwidth/7-2\tabcolsep-0.1in}
\setlength{\cellWidthe}{\columnwidth/7-2\tabcolsep-0.0in}
\setlength{\cellWidthf}{\columnwidth/7-2\tabcolsep-0.0in}
\setlength{\cellWidthg}{\columnwidth/7-2\tabcolsep-0.0in}
\scalebox{1}[1]{\begin{tabularx}{\columnwidth}{>{\PreserveBackslash\centering}m{\cellWidtha}
>{\PreserveBackslash\centering}m{\cellWidthb}
>{\PreserveBackslash\centering}m{\cellWidthc}
>{\PreserveBackslash\centering}m{\cellWidthd}
>{\PreserveBackslash\centering}m{\cellWidthe}
>{\PreserveBackslash\centering}m{\cellWidthf}
>{\PreserveBackslash\centering}m{\cellWidthg}}
\toprule
	\textbf{Factors} & \textbf{P-Wave Coeff. }& \textbf{Std. Err.} &\textbf{ t(47)} & \textbf{\emph{p}} & \boldmath{$-$}\textbf{95.\% Cnf. Limt }& \textbf{+95.\% Cnf. Limt} \\ \midrule
	       CD & 238.4 & 906.019 & 0.26314 & 0.793591 & $-$1584.3 & 2061.1 \\ \noalign{\hrule height .5pt}
       \rowcolor{green!10}  CE & 4819.4 & 906.019 & 5.31932 & 0.000003 & 2996.7 & 6642.1 \\ \noalign{\hrule height .5pt}
        DE & $-$239.9 & 906.019 & $-$0.26481 & 0.792311 & $-$2062.6 & 1582.7 \\ \noalign{\hrule height .5pt}
      \rowcolor{green!10}   ABC & $-$16,343.2 & 6380.749 & $-$2.56133 & 0.013696 & $-$29,179.6 & $-$3506.8 \\ \noalign{\hrule height .5pt}
        ABD & 4655.7 & 6380.749 & 0.72965 & 0.469228 & $-$8180.7 & 17,492.1 \\ \midrule
        ABE & 7752.7 & 6380.749 & 1.21501 & 0.230431 & $-$5083.7 & 20,589.1 \\ \midrule
        ACD & 10,057.6 & 6380.749 & 1.57624 & 0.121677 & $-$2778.8 & 22,894.0 \\ \midrule
        ACE & $-$3180.4 & 6380.749 & $-$0.49843 & 0.620506 & $-$16,016.8 & 9656.1 \\ \noalign{\hrule height .5pt}
       \rowcolor{green!10}  ADE & $-$25542.1 & 6380.749 & $-$4.00299 & 0.000221 & $-$38,378.5 & $-$12,705.7 \\ \noalign{\hrule height .5pt}
        BCD & 5030.5 & 6380.749 & 0.78839 & 0.434430 & $-$7805.9 & 17,866.9 \\ \midrule
        BCE & $-$8617.9 & 6380.749 & $-$1.35061 & 0.183289 & $-$21,454.3 & 4218.5 \\ \midrule
        BDE & 6423.5 & 6380.749 & 1.00670 & 0.319233 & $-$6412.9 & 19,259.9 \\ \midrule
        CDE & 3468.8 & 6380.749 & 0.54363 & 0.589265 & $-$9367.6 & 16,305.2 \\ \midrule
        AB(A-B) & 1873.4 & 1923.886 & 0.97375 & 0.335164 & $-$1997.0 & 5743.7 \\ \noalign{\hrule height .5pt}
      \rowcolor{green!10}  AC(A-C) & $-$7129.4 & 1923.886 & $-$3.70572 & 0.000555 & $-$10,999.7 & $-$3259.0 \\ \noalign{\hrule height .5pt}
        AD(A-D) & $-$3098.9 & 1923.886 & $-$1.61077 & 0.113926 & $-$6969.3 & 771.4 \\ \noalign{\hrule height .5pt}
      \rowcolor{green!10}   AE(A-E) & 6913.4 & 1923.886 & 3.59347 & 0.000779 & 3043.1 & 10,783.8 \\ \noalign{\hrule height .5pt}
        \rowcolor{green!10} BC(B-C) & $-$4438.2 & 1923.886 & $-$2.30690 & 0.025515 & $-$8308.6 & $-$567.8 \\ \noalign{\hrule height .5pt}
        BD(B-D) & 59.4 & 1923.886 & 0.03089 & 0.975491 & $-$3810.9 & 3929.8 \\ \noalign{\hrule height .5pt}
     \rowcolor{green!10}    BE(B-E) & 4666.2 & 1923.886 & 2.42541 & 0.019186 & 795.9 & 8536.6 \\ \noalign{\hrule height .5pt}
        CD(C-D) & 472.9 & 1923.886 & 0.24582 & 0.806890 & $-$3397.4 & 4343.3 \\ \midrule
        CE(C-E) & 2542.3 & 1923.886 & 1.32143 & 0.192754 & $-$1328.1 & 6412.6 \\ \midrule
        DE(D-E) & 818.4 & 1923.886 & 0.42541 & 0.672482 & $-$3051.9 & 4688.8 \\ \bottomrule
 \end{tabularx}}
\footnotesize{Factors indicate combination of binders. All the factors included in the cubic model. The~color-marked rows indicate a significance level of values below the \emph{p}~=~0.05. Coeffs (recoded comps); Var.:Vp\_rod\_14; R-sqr~=~0.9207; Adj: 0.8633 (Resultat\_rev091123.sta). Five-factor mixture design; Mixture total~=~1, 82 Runs. DV: Vp\_rod\_14; MS Residual~=~82,847.91. Std. Err. is a standard error (standard deviation of sample distribution).}
  \end{table}
	
	













     
Table~\ref{tab:4} shows the remaining effects on day {14}th after the backward~elimination. 


We observe here that Portland cement and lime have different effects on stabilisation (compare P-wave coefficient for cement P-wave~=~2572.1 against P-wave~=~872.6 for lime), since they differ in their chemical nature, although~both provide calcium, which is necessary for soil hardening. A~notably high value is shown by bio fly ash (P-wave~=~1192.2) due to its mode of reaction with particles. Remarkably high values are also demonstrated by the combinations AC (cement/GGBFS, P-wave~=~4811.3), AD (cement/energy fly ash P-wave~=~3085.3), BC (lime/slag P-wave~=~6901.9) and CE (GGBFS/bio fly ash, P-wave~=~4756.8). This well demonstrates that reaction products may eventually differ in real cases of soil stabilisation by blended binders. Moreover, certain triple combinations demonstrated negative values, such as, for~instance, ADE (Portland cement/energy fly ash/bio fly ash, P-waves~=~$-$29,959.3). Clearly, adding cement as a stabilising agent is beneficial for soil, as~it produces strength-developing hydration product in a soil mixture through bonding with particles, which finally leads to the improved strength of~soil. 


 \begin{table}[H]
\caption{{Remaining} significant effects on day {14} after the backward~elimination. \label{tab:4}}
\setlength{\cellWidtha}{\columnwidth/7-2\tabcolsep+0.2in}
\setlength{\cellWidthb}{\columnwidth/7-2\tabcolsep+0.0in}
\setlength{\cellWidthc}{\columnwidth/7-2\tabcolsep-0.1in}
\setlength{\cellWidthd}{\columnwidth/7-2\tabcolsep-0.1in}
\setlength{\cellWidthe}{\columnwidth/7-2\tabcolsep-0.0in}
\setlength{\cellWidthf}{\columnwidth/7-2\tabcolsep-0.0in}
\setlength{\cellWidthg}{\columnwidth/7-2\tabcolsep-0.0in}
\scalebox{1}[1]{\begin{tabularx}{\columnwidth}{>{\PreserveBackslash\centering}m{\cellWidtha}
>{\PreserveBackslash\centering}m{\cellWidthb}
>{\PreserveBackslash\centering}m{\cellWidthc}
>{\PreserveBackslash\centering}m{\cellWidthd}
>{\PreserveBackslash\centering}m{\cellWidthe}
>{\PreserveBackslash\centering}m{\cellWidthf}
>{\PreserveBackslash\centering}m{\cellWidthg}}
\toprule
        \textbf{Factors} &\textbf{ P-Wave Coeff.} & \textbf{Std. Err. }& \textbf{t(68)} & \textbf{\emph{p}} & \boldmath{$-$}\textbf{95.\% Cnf. Limt} & \textbf{+95.\% Cnf. Limt} \\ \midrule
        (A) Cement & 2572.1 & 149.090 & 17.25165 & 0.000000 & 2274.6 & 2869.6 \\ \midrule
        (B) Lime & 872.6 & 131.407 & 6.64061 & 0.000000 & 610.4 & 1134.8 \\ \midrule
        (C) GGBFS & 362.3 & 168.836 & 2.14575 & 0.035465 & 25.4 & 699.2 \\ \midrule
        (D) E\_Fly ash & 385.0 & 127.497 & 3.01993 & 0.003559 & 130.6 & 639.4 \\ \midrule
        (E) B\_Fly ash & 1192.2 & 134.910 & 8.83687 & 0.000000 & 923.0 & 1461.4 \\ \midrule
        AC & 4811.3 & 754.640 & 6.37562 & 0.000000 & 3305.4 & 6317.2 \\ \midrule
        AD & 3085.3 & 768.942 & 4.01236 & 0.000152 & 1550.9 & 4619.7 \\ \midrule
        BC & 6901.9 & 753.201 & 9.16339 & 0.000000 & 5398.9 & 8404.9 \\ \midrule
        CE & 4756.8 & 756.913 & 6.28445 & 0.000000 & 3246.4 & 6267.2 \\ \midrule
        ADE & $-$29,959.3 & 5994.247 & $-$4.99801 & 0.000004 & $-$41,920.6 & $-$17,998.0 \\ \midrule
        AC(A-C) & $-$6802.8 & 1990.624 & $-$3.41743 & 0.001071 & $-$10,775.0 & $-$2830.6 \\ \midrule
        AE(A-E) & 6630.1 & 1988.645 & 3.33396 & 0.001389 & 2661.8 & 10,598.3 \\ \midrule
        BC(B-C) & $-$4518.4 & 1984.089 & $-$2.27734 & 0.025916 & $-$8477.6 & $-$559.3 \\ \midrule
        BE(B-E) & 3976.0 & 1981.513 & 2.00657 & 0.048774 & 22.0 & 7930.1 \\ \bottomrule
\end{tabularx}}
   \footnotesize{A significance level below the \emph{p}~=~0.05 is noted for all the recorded values in Table~\ref{tab:4}. Coeffs (recoded comps); Var.:Vp\_rod\_14; R-sqr~=~0.8728; Adj: 0.8484 (Resultat\_rev091123.sta). Five-factor mixture design; Mixture total~=~1, 82 Runs. DV: Vp\_rod\_14; MS Residual~=~91,875.45. Factors are the combinations of binders.}
\end{table}







The inspection of Table~\ref{tab:5}, showing the results showing P-wave velocities on the day {28}th of curing period by temperature 46 $^\circ$C, produces following remarks. The~difference in values of the P-wave velocities, compared to the previous Tables on days {7}th and {14}th implies that the time of the reaction plays an important role on the stabilisation in soils, which generally results from the chemical processes caused by the stabilising agent as binder or the blended mixture of double of triple binders. These factors include the two different processes. The~flocculation effects come to the effect immediate upon the contact of binder with soil~particles.

On the contrary, pozzolanic and hydration effects are generally time-dependent. Therefore, we can see the differences in stabilisation, as~reflected by the P-waves speed on days {7}th, {14}th and {28}th by the same stabilising agents. For~instance, the combinations of cement/lime (blend AB) or cement/GGBSF (blend AC) for the day {28}th, which give the P-wave coefficients as 928.1 and 5917.3 against similar values on day {14}th and {7}th. Thus, P-wave coefficients for the combination AB for day {14}th:  828.6 and for the day 7: 730.2, respectively. Likewise, for~the combination AC P-waves are recorded as 1345.3 for day {7}th and 5229.2 for day {14}th. This can be explained by the inherent nature of the Portland cement, in~addition to the added binders, which contributes to the increase of strength and plasticity reduction over the time of~curing. 

 \begin{table}[H]
 \renewcommand{\arraystretch}{1.3}
 \tablesize{\footnotesize}
\caption{The effects from various binders on soil stabilisation on day {28}th of curing, measured by P-wave~velocities.}\label{tab:5}
\newcolumntype{C}{>{\centering\arraybackslash}X}
\begin{tabularx}{\textwidth}{CCCCCCC}
\toprule
	     \textbf{Factors} &\textbf{ P-Wave Coeff.} & \textbf{Std. Err. }& \textbf{t(47)} & \textbf{\emph{p}} & \boldmath{$-$}\textbf{95.\% Cnf. Limt} & \textbf{+95.\% Cnf. Limt} \\ \noalign{\hrule height .5pt}
       \rowcolor{green!10}  (A) Cement & 2757.6 & 227.599 & 12.11585 & 0.000000 & 2299.7 & 3215.4 \\ \noalign{\hrule height .5pt}
      \rowcolor{green!10}   (B) Lime & 1094.6 & 227.599 & 4.80915 & 0.000016 & 636.7 & 1552.4 \\ \noalign{\hrule height .5pt}
        (C) GGBFS & 353.0 & 227.599 & 1.55109 & 0.127589 & $-$104.8 & 810.9 \\ \noalign{\hrule height .5pt}
        (D) E\_Fly ash & 213.6 & 227.599 & 0.93831 & 0.352883 & $-$244.3 & 671.4 \\ \noalign{\hrule height .5pt}
       \rowcolor{green!10}  (E) B\_Fly ash & 1798.2 & 227.599 & 7.90068 & 0.000000 & 1340.3 & 2256.1 \\ \noalign{\hrule height .5pt}
        AB & 928.1 & 1018.836 & 0.91097 & 0.366960 & $-$1121.5 & 2977.8 \\ \noalign{\hrule height .5pt}
     \rowcolor{green!10}    AC & 5917.3 & 1018.836 & 5.80793 & 0.000001 & 3867.7 & 7967.0 \\ \noalign{\hrule height .5pt}
       \rowcolor{green!10}  AD & 2926.6 & 1018.836 & 2.87252 & 0.006093 & 877.0 & 4976.3 \\ \noalign{\hrule height .5pt}
        AE & $-$1973.4 & 1018.836 & $-$1.93691 & 0.058781 & $-$4023.0 & 76.2 \\ \noalign{\hrule height .5pt}
       \rowcolor{green!10}  BC & 8789.8 & 1018.836 & 8.62726 & 0.000000 & 6740.1 & 10,839.4 \\ \noalign{\hrule height .5pt}
        BD & 1194.6 & 1018.836 & 1.17251 & 0.246903 & $-$855.0 & 3244.2 \\ \noalign{\hrule height .5pt}
        BE & $-$1261.8 & 1018.836 & $-$1.23850 & 0.221680 & $-$3311.5 & 787.8 \\ \noalign{\hrule height .5pt}
        CD & 122.4 & 1018.836 & 0.12018 & 0.904855 & $-$1927.2 & 2172.1 \\ \noalign{\hrule height .5pt}
      \rowcolor{green!10}   CE & 6434.7 & 1018.836 & 6.31573 & 0.000000 & 4385.1 & 8484.3 \\ \noalign{\hrule height 1pt}       
           DE & 388.6 & 1018.836 & 0.38137 & 0.704650 & $-$1661.1 & 2438.2 \\ \noalign{\hrule height .5pt}
      \rowcolor{green!10}   ABC & $-$14,519.3 & 7175.280 & $-$2.02352 & 0.048727 & $-$28,954.1 & $-$84.5 \\ \noalign{\hrule height .5pt}
        ABD & 3311.1 & 7175.280 & 0.46146 & 0.646600 & $-$11,123.7 & 17,745.9 \\ \noalign{\hrule height .5pt}
        ABE & 9603.0 & 7175.280 & 1.33834 & 0.187224 & $-$4831.8 & 24,037.8 \\ \noalign{\hrule height .5pt}
        ACD & 13,143.3 & 7175.280 & 1.83174 & 0.073332 & $-$1291.5 & 27,578.1 \\ \noalign{\hrule height .5pt}
        ACE & $-$6703.8 & 7175.280 & $-$0.93429 & 0.354928 & $-$21,138.6 & 7731.0 \\ \noalign{\hrule height .5pt}
      \rowcolor{green!10}   ADE & $-$26,777.2 & 7175.280 & $-$3.73187 & 0.000512 & $-$41,212.0 & $-$12,342.4 \\ \noalign{\hrule height .5pt}
        BCD & 12,947.4 & 7175.280 & 1.80444 & 0.077572 & $-$1487.4 & 27,382.2 \\ \noalign{\hrule height .5pt}
        BCE & $-$11,210.9 & 7175.280 & $-$1.56243 & 0.124896 & $-$25,645.7 & 3223.9 \\ \noalign{\hrule height .5pt}
        BDE & 6493.1 & 7175.280 & 0.90493 & 0.370118 & $-$7941.7 & 20,927.9 \\ \noalign{\hrule height .5pt}
        CDE & 9002.6 & 7175.280 & 1.25466 & 0.215805 & $-$5432.2 & 23,437.4 \\ \noalign{\hrule height .5pt}
        AB(A-B) & 1405.5 & 2163.448 & 0.64968 & 0.519063 & $-$2946.8 & 5757.8 \\ \noalign{\hrule height .5pt}
    \rowcolor{green!10}     AC(A-C) & $-$7717.7 & 2163.448 & $-$3.56730 & 0.000843 & $-$12,070.0 & $-$3365.4 \\ \noalign{\hrule height .5pt}
        AD(A-D) & $-$3623.9 & 2163.448 & $-$1.67507 & 0.100562 & $-$7976.2 & 728.4 \\ \noalign{\hrule height .5pt}
    \rowcolor{green!10}     AE(A-E) & 7851.3 & 2163.448 & 3.62908 & 0.000700 & 3499.0 & 12,203.6 \\ \noalign{\hrule height .5pt}
     \rowcolor{green!10}    BC(B-C) & $-$4806.3 & 2163.448 & $-$2.22158 & 0.031165 & $-$9158.6 & $-$454.0 \\ \noalign{\hrule height .5pt}
        BD(B-D) & $-$678.2 & 2163.448 & $-$0.31346 & 0.755320 & $-$5030.4 & 3674.1 \\ \noalign{\hrule height .5pt}
       \rowcolor{green!10}  BE(B-E) & 4889.2 & 2163.448 & 2.25992 & 0.028501 & 536.9 & 9241.5 \\ \noalign{\hrule height .5pt}
        CD(C-D) & $-$202.8 & 2163.448 & $-$0.09372 & 0.925729 & $-$4555.1 & 4149.5 \\ \noalign{\hrule height .5pt}
    \rowcolor{green!10}     CE(C-E) & 6361.9 & 2163.448 & 2.94064 & 0.005068 & 2009.6 & 10,714.2 \\ \noalign{\hrule height .5pt}
        DE(D-E) & $-$1936.1 & 2163.448 & $-$0.89490 & 0.375400 & $-$6288.4 & 2416.2 \\ \bottomrule
\end{tabularx}
\footnotesize{Factors signify the combination of binders. All the factors included in the cubic model. The~color-marked rows indicate a significance level of values below the \emph{p}~=~0.05. Coeffs (recoded comps); Var.:Vp\_rod\_28\_f; R-sqr~=~0.9152; Adj: 0.8538 (Resultat\_rev091123.sta) Five-factor mixture design; Mixture total~=~1, 82 Runs. DV: Vp\_rod\_28\_f; MS Residual~=~104,764.9. Std. Err. is a standard error (standard deviation of sample distribution).}
\end{table}


\newpage
Table~\ref{tab:6} shows the remaining significant effects on soil stabilisation from different binders on day {28}th after the backward elimination in the regression analysis. The~dynamics in P-waves which generally increases over time is explained by the increased soil strength as a result of the stabilisation, where plasticity of soil is reduced and particles are bonded to each other in a soil-blend mixture. Thus, the~compressive strength and load-bearing properties of soil are improved, which results in the increased stiffness and compaction of soil and reduced porosity. This corresponds to the higher of the P-wave velocity which rises in a more dense structure, due to the chemical processes in soil caused by binders. 

\begin{table}[H]
\tablesize{\footnotesize}
\caption{{Remaining} significant effects on day {28}th after the backward~elimination. \label{tab:6}}
\setlength{\cellWidtha}{\columnwidth/7-2\tabcolsep+0.2in}
\setlength{\cellWidthb}{\columnwidth/7-2\tabcolsep+0.0in}
\setlength{\cellWidthc}{\columnwidth/7-2\tabcolsep-0.1in}
\setlength{\cellWidthd}{\columnwidth/7-2\tabcolsep-0.1in}
\setlength{\cellWidthe}{\columnwidth/7-2\tabcolsep-0.0in}
\setlength{\cellWidthf}{\columnwidth/7-2\tabcolsep-0.0in}
\setlength{\cellWidthg}{\columnwidth/7-2\tabcolsep-0.0in}
\scalebox{1}[1]{\begin{tabularx}{\columnwidth}{>{\PreserveBackslash\centering}m{\cellWidtha}
>{\PreserveBackslash\centering}m{\cellWidthb}
>{\PreserveBackslash\centering}m{\cellWidthc}
>{\PreserveBackslash\centering}m{\cellWidthd}
>{\PreserveBackslash\centering}m{\cellWidthe}
>{\PreserveBackslash\centering}m{\cellWidthf}
>{\PreserveBackslash\centering}m{\cellWidthg}}
\toprule
           \textbf{Factors} &\textbf{ P-Wave Coeff.} & \textbf{Std. Err. }& \textbf{t(69)} & \textbf{\emph{p}} & \boldmath{$-$}\textbf{95.\% Cnf. Limt} & \textbf{+95.\% Cnf. Limt} \\ \midrule
        (A) Cement & 2641.7 & 188.071 & 14.04626 & 0.000000 & 2266.5 & 3016.9 \\ \midrule
       (B) Lime & 1234.2 & 159.337 & 7.74602 & 0.000000 & 916.4 & 1552.1 \\ \midrule
       (C) GGBFS & 535.2 & 213.044 & 2.51231 & 0.014337 & 110.2 & 960.2 \\ \midrule
       (D) E\_Fly ash & 559.9 & 160.875 & 3.48025 & 0.000873 & 238.9 & 880.8 \\ \midrule
       (E) B\_Fly ash & 1504.3 & 169.277 & 8.88679 & 0.000000 & 1166.6 & 1842.0 \\ \midrule
       AC & 5570.6 & 951.996 & 5.85146 & 0.000000 & 3671.4 & 7469.8 \\ \midrule
       AD & 3190.2 & 970.228 & 3.28813 & 0.001589 & 1254.7 & 5125.8 \\ \midrule
        BC & 7892.0 & 951.349 & 8.29563 & 0.000000 & 5994.2 & 9789.9 \\ \midrule
        CE & 6318.7 & 950.481 & 6.64791 & 0.000000 & 4422.6 & 8214.9 \\ \midrule
       ADE & $-$29,928.3 & 7562.989 & $-$3.95721 & 0.000182 & $-$45,016.1 & $-$14,840.6 \\ \midrule
        AC(A-C) & $-$7511.2 & 2511.772 & $-$2.99042 & 0.003860 & $-$12,522.1 & $-$2500.4 \\ \midrule
       AE(A-E) & 7551.3 & 2509.255 & 3.00936 & 0.003653 & 2545.4 & 12,557.1 \\ \midrule
        CE(C-E) & 5855.4 & 2503.606 & 2.33880 & 0.022249 & 860.9 & 10,850.0 \\ \bottomrule
   \end{tabularx}}
     \footnotesize{A~significance level below the \emph{p}~=~0.05 is noted for all the recorded values in Table~\ref{tab:6}. Coeffs (recoded comps); Var.:Vp\_rod\_28\_f; R-sqr~=~0.8261; Adj: 0.7958 (Resultat\_rev091123.sta). Five-factor mixture design; Mixture total~=~1, 82 Runs. DV: Vp\_rod\_28\_f; MS Residual~=~146,270.9.} 
\end{table}




Table~\ref{tab:7} shows that slag is the main factor that contributes most to the compressive strength among the other binders followed by the OPC, which comes second in terms of strength. The~combination of slag and energy ash has the least impact on the compressive strength of soil. Furthermore, Table~\ref{tab:7} shows linear effects between cement, lime and slag as factors affecting soil stabilisation and gain of strength when comparing the values of P-waves, corresponding to the gain of strength over curing time of 3 months. By~closer examination of Table~\ref{tab:7}, the~effects from lime and slag that are reflected in seismic measurements (7th, 14th and 28th daily values in Table~\ref{tab:1}, Table \ref{tab:3} and Table \ref{tab:5}, respectively) cannot be confirmed for the values of strength recorded on day {90}. Here high values are demonstrated by the blend lime/slag (BC, P-wave~=~6004.0), lime/energy fly slag (BD, P-waves~=~4620.5) and cement/slag (AC, 5147.8). Table~\ref{tab:8} shows the remaining significant effects on day {90} after backward elimination in regression analysis.


\begin{table}[H]
\renewcommand{\arraystretch}{1.3}
\caption{{The effects} from various binders on soil stabilisation on day {90}th of curing time, measured by P-wave~velocities.}\label{tab:7}
\setlength{\cellWidtha}{\columnwidth/7-2\tabcolsep+0.0in}
\setlength{\cellWidthb}{\columnwidth/7-2\tabcolsep+0.0in}
\setlength{\cellWidthc}{\columnwidth/7-2\tabcolsep-0.0in}
\setlength{\cellWidthd}{\columnwidth/7-2\tabcolsep-0.0in}
\setlength{\cellWidthe}{\columnwidth/7-2\tabcolsep-0.0in}
\setlength{\cellWidthf}{\columnwidth/7-2\tabcolsep-0.0in}
\setlength{\cellWidthg}{\columnwidth/7-2\tabcolsep-0.0in}
\scalebox{1}[1]{\begin{tabularx}{\columnwidth}{>{\PreserveBackslash\centering}m{\cellWidtha}
>{\PreserveBackslash\centering}m{\cellWidthb}
>{\PreserveBackslash\centering}m{\cellWidthc}
>{\PreserveBackslash\centering}m{\cellWidthd}
>{\PreserveBackslash\centering}m{\cellWidthe}
>{\PreserveBackslash\centering}m{\cellWidthf}
>{\PreserveBackslash\centering}m{\cellWidthg}}
\noalign{\hrule height 1pt}
	\rowcolor{green!10}   \textbf{Factors}  &\textbf{ P-Wave Coeff.} & \textbf{Std. Err. }& \textbf{t(47)} & \textbf{\emph{p}} & \boldmath{$-$}\textbf{95.\% Cnf. Limt} & \textbf{+95.\% Cnf. Limt} \\ \noalign{\hrule height .5pt}
  \rowcolor{green!10}       (A) Cement & 5992.1 & 338.88 & 17.68189 & 0.000000 & 5310.3 & 6673.8 \\ \noalign{\hrule height .5pt}
   \rowcolor{green!10}     (B) Lime & 1387.2 & 338.88 & 4.09355 & 0.000166 & 705.5 & 2069.0 \\ \noalign{\hrule height .5pt}
    \rowcolor{green!10}     (C) GGBFS & 6910.7 & 338.88 & 20.39270 & 0.000000 & 6229.0 & 7592.4 \\ \noalign{\hrule height .5pt}
        (D) E\_Fly ash & 263.6 & 338.88 & 0.77791 & 0.440522 & $-$418.1 & 945.4 \\ \noalign{\hrule height .5pt}
    \rowcolor{green!10}     (E) B\_Fly ash & 1577.6 & 338.88 & 4.65543 & 0.000027 & 895.9 & 2259.4 \\ \noalign{\hrule height .5pt}
       AB & $-$2448.3 & 1516.98 & $-$1.61394 & 0.113235 & $-$5500.1 & 603.5 \\ \noalign{\hrule height .5pt}
   \rowcolor{green!10}      AC & 5147.8 & 1516.98 & 3.39344 & 0.001410 & 2096.0 & 8199.6 \\ \noalign{\hrule height .5pt}
        AD & 1775.3 & 1516.98 & 1.17030 & 0.247780 & $-$1276.4 & 4827.1 \\ \noalign{\hrule height .5pt}
       AE & 1710.6 & 1516.98 & 1.12761 & 0.265208 & $-$1341.2 & 4762.3 \\ \noalign{\hrule height .5pt}
   \rowcolor{green!10}      BC & 6004.0 & 1516.98 & 3.95785 & 0.000254 & 2952.2 & 9055.8 \\ \noalign{\hrule height .5pt}
    \rowcolor{green!10}     BD & 4620.5 & 1516.98 & 3.04586 & 0.003797 & 1568.7 & 7672.3 \\ \noalign{\hrule height .5pt}
        BE & 834.3 & 1516.98 & 0.54994 & 0.584961 & $-$2217.5 & 3886.0 \\ \noalign{\hrule height .5pt}
    \rowcolor{green!10}     CD & $-$12,683.0 & 1516.98 & $-$8.36066 & 0.000000 & $-$15,734.8 & $-$9631.2 \\ \noalign{\hrule height .5pt}
      \rowcolor{green!10}   CE & 6310.9 & 1516.98 & 4.16015 & 0.000134 & 3259.1 & 9362.7 \\ \noalign{\hrule height .5pt}
        DE & 2361.3 & 1516.98 & 1.55657 & 0.126282 & $-$690.5 & 5413.1 \\ \noalign{\hrule height .5pt}
        ABC & $-$14620.6 & 10,683.55 & $-$1.36851 & 0.177658 & $-$36,113.1 & 6872.0 \\ \noalign{\hrule height .5pt}
        ABD & $-$1054.7 & 10,683.55 & $-$0.09872 & 0.921779 & $-$22,547.2 & 20,437.8 \\ \noalign{\hrule height .5pt}
       ABE & $-$6205.2 & 10,683.55 & $-$0.58082 & 0.564140 & $-$27,697.7 & 15,287.4 \\ \noalign{\hrule height .5pt}
      \rowcolor{green!10}   ACD & 45,351.2 & 10,683.55 & 4.24496 & 0.000102 & 23,858.7 & 66,843.7 \\ \noalign{\hrule height .5pt}
      \rowcolor{green!10}  ACE & $-$45,649.2 & 10,683.55 & $-$4.27285 & 0.000093 & $-$67,141.7 & $-$24,156.6 \\ \noalign{\hrule height .5pt}
        ADE & 10,067.1 & 10,683.55 & 0.94230 & 0.350859 & $-$11,425.4 & 31,559.6 \\ \noalign{\hrule height .5pt}
       \rowcolor{green!10}  BCD & 54,712.7 & 10,683.55 & 5.12122 & 0.000006 & 33,220.2 & 76,205.3 \\ \noalign{\hrule height .5pt}
        BCE & $-$21,262.1 & 10,683.55 & $-$1.99017 & 0.052406 & $-$42,754.6 & 230.4 \\ \noalign{\hrule height .5pt}
        BDE & 7008.6 & 10,683.55 & 0.65602 & 0.515009 & $-$14,483.9 & 28,501.2 \\ \noalign{\hrule height .5pt}
    \rowcolor{green!10}    CDE & 43,775.5 & 10,683.55 & 4.09747 & 0.000164 & 22,283.0 & 65,268.0 \\ \noalign{\hrule height .5pt}
        AB(A-B) & 1876.0 & 3221.24 & 0.58237 & 0.563099 & $-$4604.3 & 8356.3 \\ \noalign{\hrule height .5pt}
        AC(A-C) & $-$6009.6 & 3221.24 & $-$1.86562 & 0.068343 & $-$12,489.9 & 470.7 \\ \noalign{\hrule height .5pt}
        AD(A-D) & 648.5 & 3221.24 & 0.20132 & 0.841319 & $-$5831.8 & 7128.8 \\ \noalign{\hrule height .5pt}
        AE(A-E) & 5207.7 & 3221.24 & 1.61667 & 0.112643 & $-$1272.6 & 11,688.0 \\ \noalign{\hrule height .5pt}
    \rowcolor{green!10}     BC(B-C) & $-$8918.3 & 3221.24 & $-$2.76860 & 0.008031 & $-$15,398.6 & $-$2438.0 \\ \noalign{\hrule height .5pt}
        BD(B-D) & 419.5 & 3221.24 & 0.13024 & 0.896934 & $-$6060.8 & 6899.8 \\ \noalign{\hrule height .5pt}
        BE(B-E) & $-$887.0 & 3221.24 & $-$0.27537 & 0.784238 & $-$7367.3 & 5593.3 \\ \noalign{\hrule height .5pt}
    \rowcolor{green!10}     CD(C-D) & $-$10,632.0 & 3221.24 & $-$3.30060 & 0.001847 & $-$17,112.3 & $-$4151.7 \\ \noalign{\hrule height .5pt}
   \rowcolor{green!10}      CE(C-E) & 7909.8 & 3221.24 & 2.45551 & 0.017822 & 1429.5 & 14,390.1 \\ \noalign{\hrule height .5pt}
        DE(D-E) & $-$3409.2 & 3221.24 & $-$1.05835 & 0.295307 & $-$9889.5 & 3071.1 \\ \bottomrule
          \end{tabularx}}
    \footnotesize{All the factors included in the cubic model. The~color-marked values idicate a significance level below \emph{p}~=~0.05. Coeffs (recoded comps); Var.:UCS; R-sqr~=~0.9678; Adj: 0.9445 (Resultat\_rev091123.sta). Five-factor mixture design; Mixture total~=~1, 82 Runs. DV: UCS; MS Residual~=~232,257.1 Std. Err. is a standard error (standard deviation of sample distribution).}
\end{table}


\begin{table}[H]
\caption{{Remaining} significant effects on day {90}th after backward~elimination.  \label{tab:8}}
\setlength{\cellWidtha}{\columnwidth/7-2\tabcolsep+0.2in}
\setlength{\cellWidthb}{\columnwidth/7-2\tabcolsep+0.0in}
\setlength{\cellWidthc}{\columnwidth/7-2\tabcolsep-0.1in}
\setlength{\cellWidthd}{\columnwidth/7-2\tabcolsep-0.1in}
\setlength{\cellWidthe}{\columnwidth/7-2\tabcolsep-0.0in}
\setlength{\cellWidthf}{\columnwidth/7-2\tabcolsep-0.0in}
\setlength{\cellWidthg}{\columnwidth/7-2\tabcolsep-0.0in}
\scalebox{1}[1]{\begin{tabularx}{\columnwidth}{>{\PreserveBackslash\centering}m{\cellWidtha}
>{\PreserveBackslash\centering}m{\cellWidthb}
>{\PreserveBackslash\centering}m{\cellWidthc}
>{\PreserveBackslash\centering}m{\cellWidthd}
>{\PreserveBackslash\centering}m{\cellWidthe}
>{\PreserveBackslash\centering}m{\cellWidthf}
>{\PreserveBackslash\centering}m{\cellWidthg}}
\toprule
       \textbf{Factors} &\textbf{P-Wave Coeff.} & \textbf{Std. Err. }& \textbf{t(71)} & \textbf{\emph{p}} & \boldmath{$-$}\textbf{95.\% Cnf. Limt} & \textbf{+95.\% Cnf. Limt}\\ \midrule
        (A) Cement & 6042.5 & 238.21 & 25.3667 & 0.000000 & 5567.5 & 6517.5 \\ \midrule
        (B) Lime & 1497.8 & 235.56 & 6.3587 & 0.000000 & 1028.1 & 1967.5 \\ \midrule
        (C) GGBFS & 8499.5 & 269.19 & 31.5738 & 0.000000 & 7962.7 & 9036.3 \\ \midrule
        (D) E\_Fly ash & 1006.3 & 265.81 & 3.7857 & 0.000317 & 476.3 & 1536.3 \\ \midrule
        (E) B\_Fly ash & 1985.8 & 238.21 & 8.3363 & 0.000000 & 1510.8 & 2460.7 \\ \midrule
        CD & $-$17,791.7 & 1675.96 & $-$10.6158 & 0.000000 & $-$21,133.5 & $-$14,449.9 \\ \midrule
        ACD & 60,200.0 & 12,482.71 & 4.8227 & 0.000008 & 35,310.2 & 85,089.8 \\ \midrule
       ACE & $-$24,816.7 & 12,149.73 & $-$2.0426 & 0.044809 & $-$49,042.6 & $-$590.9 \\ \midrule        
       BCD & 78,876.2 & 12,501.89 &6.3091 & 0.000000 & 53,948.1 & 103,804.2 \\ \midrule
        CDE & 59,769.4 & 12,482.71 & 4.7882 & 0.000009 & 34,879.6 & 84,659.2 \\ \midrule
        CD(C-D) & $-$11,168.5 & 4107.77 & $-$2.7189 & 0.008227 & $-$19,359.2 & $-$2977.9 \\ \bottomrule
   \end{tabularx}}
     \footnotesize{A~significance level \emph{p}~=~0.05 is noted for all the recorded values. Coeffs (recoded comps); Var.:UCS; R-sqr~=~0.9164; Adj: 0.9046 (Resultat\_rev091123.sta). Five-factor mixture design; Mixture total~=~1, 82 Runs. DV: UCS; MS Residual~=~39,9436.4.}
\end{table}






The combination of lime and slag has a significant effect on hardening process of soil, but~no effects on  final strength development. These observations also confirm previous study \cite{Lindh2001} on the three different types of the fine-grained moraines, of~which the two samples with the largest clay content demonstrated a significant interaction between lime and slag, while the third sample with the lowest clay content shown the same pattern, as shown in Table~\ref{tab:7}. This indicates that bearing gravel used in the study is inert and therefore well suited to study the effects of different binders and their interactions. The~remaining significant effects on day {90}th after backward elimination show notable values for single binders slag (P-wave~=~8499.5) and cement (P-wave~=~6042.5), as~well as triple combination of blended mixture of binders Portland cement/ACD/slag GGBFS/fly ash (P-wave~{=~60,200.0}), as~shown in Table \ref{tab:8}. 

%%%%%%%%%%%%%%%%%%%%%%%%%%%%%%%%%%%%%%%%%%
\section{Conclusions}
In this study, we demonstrated that blended binders perform better on soil stabilisation, as~compared to the single binders. The~experiments for optimal selection of binder combinations and amounts reduce the financial risks of the project, which are possible if soil is not sufficiently stabilised, and minimise the environmental impact through the reduced amounts of cement by replacement from other binders. Because~the costs of the laboratory tests are expensive, optimising blended binders is essential for high performance of soil in real  construction projects. In~turn, the~decisions based on the results of soil stabilisation and evaluated soil strength support solutions for technical, economic and environmental processes. Trial planning also provides the support for evaluation if optimal mixture is within a wide or narrow mixture range of blended binders. If~it is narrow, then greater precision is required in the decisions, compared to the cases with a wide range. In~contrast, if~the range of binder choices is wide, other parameters can be decisive for final solutions. These includes, for~instance, geotechnical properties of soil, economics or delivery~security.

Stabilisation of poor subgrade soil intended for road construction is aimed at improving its engineering properties, commonly estimated as strength and stiffness. The~growth of strength and stiffness is strongly dependent on time of curing and the effects from the stabilising agents which can be used both as single binders and as blended mixtures. Using mixed binders as stabilising agents may have diverse effects on the properties of soil. Therefore, the~experimental testing of the effects of blenders is required. In~this study, the comparative analysis was performed to assess the functional behavior of individual and blended (binary and triple) binders and curing time of 90 days using P-wave velocity. The~development in P-wave velocity was compared both for individual and blended binders over the complete period of curing. The results indicate that strength growth in stabilised soil samples is nonlinear in both time and content of~binders. 

The analysis of the effects from various agents demonstrated that GGBFS contributes the most to the compressive strength development, followed by OPC. Moreover, we noted that blended mixtures perform better than single binders. Thus, the~combination of lime and slag has a notable effect on soil hardening and final strength development. The~combination of slag and energy fly ash contributes to the gain of strength, but~has the least impact on compared to other blends. Factor test combining blended binders of various mixtures, enabled to investigate, under~which circumstances two factors interact during the process of soil stabilisation, so that the quality of the final product to be approved or not, e.g.,~if soil treatment takes place in the sub-area with a high contaminant content and high water content. Tailored prevention strategies in road construction can be developed {based on this procedure, by~selectively increasing the amount of binders in the areas that constitute to the risk zones or regulating the permitted water ratio in the soil base material. From~the geotechnical and engineering perspectives in context of road or highway construction, soil stabilisation should ideally have a lifespan of 80 years or more. From~an environmental perspective, in~case of contaminated soil, stabilisation of specimens should have a risk-reducing effect over the significantly longer periods of~time. 

In order to ensure that final product meets the approved requirements, both geotechnical construction of roads and physiochemical improvements of soil quality should be considered. At~the same time, base material exhibits varying properties within the treatment surface, therefore the robustness of the binder mixture should be tested by factor test, as~demonstrated in this study. Using simplex design method of binder combinations we estimated and compared the effects from blended mixtures and single binders (%MDPI: please add right bracket or we will remove this. -- Answer: added, corrected.
OPC, lime, GGBFS type Merit 5000, energy fly ash and bio flay ash} on soil stabilisation.

The robustness of soil stabilisation means that final binder mixture is able to provide the approved geotechnical and environmental quality of soil intended for road construction. This is necessary both for the short-term and long-term period of curing under all the conditions prevailing in soil treatment. The~procedure of stabilisation can lead to the unnecessarily high binder costs and increase climate emissions through the binder's production chain (especially for OPC) or the unnecessarily high technical performance given the basic requirements that govern the quality of the final product. In~the traditional tests that consider only one factor at a time, the~robustness of blending binder is often handled by the increased amount of binder, so that the results of the provided experiment ensure the avoided uncertainties that may exist: the individual physiochemical properties of soil and local environmental conditions of real project, e.g.,~specifics of roads and local climate. In this paper, we designed adaptation algorithms for selecting optimal binder combinations using simplex design combinatorics to improve soil properties and geotechnical parameters through stabilisation for industrial projects of road and highway construction.

%%%%%%%%%%%%%%%%%%%%%%%%%%%%%%%%%%%%%%%%%%
\vspace{6pt} 

\authorcontributions{\textls[-35]{Conceptualisation, methodology, software, validation, visualisation, supervision---P.L. (Per Lindh); formal analysis, investigation, resources, data curation, project administration, funding acquisition---P.L. (Polina Lemenkova); writing---original draft preparation, P.L. (Polina Lemenkova); writing---review and editing, P.L. (Per Lindh) and P.L. (Polina Lemenkova). All authors have read and agreed to the published version of the~manuscript.}}

\funding{This %MDPI: Please carefully check the funding information especially the grant number.  Please note that Funding information can not be changed after publication. -- Answer: yes, funding information is correct.
 research was funded by Materials Editorial Office of the MDPI, by~providing 100\% APC fee waiver for publication of this manuscript.}

\institutionalreview{Not applicable.}

\informedconsent{Not applicable.}

\dataavailability{Not applicable.} 

\acknowledgments{The authors are thankful to the three anonymous reviewers for their comments and remarks which improved both text and graphics in the initial version of the~manuscript.}

\conflictsofinterest{The authors declare no conflict of~interest.} 


\abbreviations{Abbreviations}{
The following abbreviations are used in this manuscript:\\

\noindent 
\begin{tabular}{@{}ll}
CHP & Combined Heat and Power \\
GGBFS & Ground Granulated Blast Furnace Slag \\
ICP & Integrated Circuit Piezoelectric \\
ISO & International Organization for Standardization \\
OPC & Ordinary Portland Cement \\
PCB & PicoCoulomB \\
SGI & Swedish Geotechnical Institute \\
SIS & Swedish Institute for Standards \\
UCS & Unconfined/Uniaxial Compressive Strength \\
USCS & Unified Soil Classification System \\
\end{tabular}
}
%%%%%%%%%%%%%%%%%%%%%%%%%%%%%%%%%%%%%%%%%%


\appendixtitles{no} % Leave argument "no" if all appendix headings stay EMPTY (then no dot is printed after "Appendix A"). If~the appendix sections contain a heading then change the argument to "yes".
\appendixstart
\appendix
\section[\appendixname~\thesection]{}

%------------------------>



\begin{table}[H]
\renewcommand{\arraystretch}{1.3}
\tablesize{\scriptsize}
\caption{Schedule %MDPI: Missing citation for Appendix A. Appendix A has not been cited in the main text. Please add the citation of Appendix A or table A1. -- Answer: citation added in the phrase "Detailed schedule of the binder recipes for each soil specimen is presented in Table \ref{tabA1}. " (\subsubsection{2.1.2. Binders})
 of the recipes for each soil specimen. Replicate 2 is a double~attempt. \label{tabA1}}
\begin{adjustwidth}{-\extralength}{0cm}
		\newcolumntype{C}{>{\centering\arraybackslash}X}
		\begin{tabularx}{\fulllength}{CCCCCCCCCCCCCC}
			\noalign{\hrule height 1pt} 
			\rowcolor{green!10} 
\textbf{PN} & \textbf{R} & \textbf{Cement} & \textbf{Lime} & \textbf{Slag} & \textbf{Flyash1} & \textbf{Flyash2} & \textbf{PN} & \textbf{R} & \textbf{Cement} & \textbf{Lime} & \textbf{Slag} & \textbf{Flyash1} & \textbf{Flyash2}\\ \noalign{\hrule height .5pt} 
39 & 1 & 0.100000 & 0.100000 & 0.100000 & 0.600000 & 0.100000 &   28 & 1 & 0.333333 & 0.333333 & 0.000000 & 0.000000 & 0.333333 \\ \midrule
        65 & 2 & 0.000000 & 0.000000 & 0.666667 & 0.000000 & 0.333333 &   29 & 1 & 0.333333 & 0.000000 & 0.333333 & 0.333333 & 0.000000\\ \midrule
        82 & 2 & 0.200000 & 0.200000 & 0.200000 & 0.200000 & 0.200000 &   71 & 2 & 0.333333 & 0.000000 & 0.333333 & 0.000000 & 0.333333 \\ \midrule
        37 & 1 & 0.100000 & 0.600000 & 0.100000 & 0.100000 & 0.100000 & 33 & 1 & 0.000000 & 0.333333 & 0.333333 & 0.000000 & 0.333333\\ \midrule
        49 & 2 & 0.333333 & 0.000000 & 0.000000 & 0.666667 & 0.000000 & 11 & 1 & 0.000000 & 0.333333 & 0.000000 & 0.666667 & 0.000000\\ \midrule
        35 & 1 & 0.000000 & 0.000000 & 0.333333 & 0.333333 & 0.333333 &  79 & 2 & 0.100000 & 0.100000 & 0.600000 & 0.100000 & 0.100000 \\ \midrule
        73 & 2 & 0.000000 & 0.333333 & 0.333333 & 0.333333 & 0.000000 & 13 & 1 & 0.000000 & 0.000000 & 0.333333 & 0.666667 & 0.000000 \\ \midrule
        68 & 2 & 0.333333 & 0.333333 & 0.000000 & 0.333333 & 0.000000 & 61 & 2 & 0.000000 & 0.666667 & 0.333333 & 0.000000 & 0.000000\\ \midrule
        76 & 2 & 0.000000 & 0.000000 & 0.333333 & 0.333333 & 0.333333 & 19 & 1 & 0.666667 & 0.000000 & 0.000000 & 0.000000 & 0.333333\\ \midrule
        16 & 1 & 0.666667 & 0.333333 & 0.000000 & 0.000000 & 0.000000 & 17 & 1 & 0.666667 & 0.000000 & 0.333333 & 0.000000 & 0.000000 \\ \midrule
        15 & 1 & 0.000000 & 0.000000 & 0.000000 & 0.333333 & 0.666667 &  63 & 2 & 0.000000 & 0.666667 & 0.000000 & 0.000000 & 0.333333\\ \midrule
        41 & 1 & 0.200000 & 0.200000 & 0.200000 & 0.200000 & 0.200000 & 81 & 2 & 0.100000 & 0.100000 & 0.100000 & 0.100000 & 0.600000\\ \midrule
        66 & 2 & 0.000000 & 0.000000 & 0.000000 & 0.666667 & 0.333333 & 31 & 1 & 0.333333 & 0.000000 & 0.000000 & 0.333333 & 0.333333\\ \midrule
        67 & 2 & 0.333333 & 0.333333 & 0.333333 & 0.000000 & 0.000000 &  44 & 2 & 0.000000 & 0.000000 & 1.000000 & 0.000000 & 0.000000\\ \midrule
        42 & 2 & 1.000000 & 0.000000 & 0.000000 & 0.000000 & 0.000000 &  53 & 2 & 0.000000 & 0.333333 & 0.000000 & 0.000000 & 0.666667\\ \midrule
        50 & 2 & 0.333333 & 0.000000 & 0.000000 & 0.000000 & 0.666667 &  2 & 1 & 0.000000 & 1.000000 & 0.000000 & 0.000000 & 0.000000\\ \midrule
        38 & 1 & 0.100000 & 0.100000 & 0.600000 & 0.100000 & 0.100000 &   3 & 1 & 0.000000 & 0.000000 & 1.000000 & 0.000000 & 0.000000\\ \midrule
        6 & 1 & 0.333333 & 0.666667 & 0.000000 & 0.000000 & 0.000000 & 9 & 1 & 0.333333 & 0.000000 & 0.000000 & 0.000000 & 0.666667\\ \midrule
        69 & 2 & 0.333333 & 0.333333 & 0.000000 & 0.000000 & 0.333333 &   46 & 2 & 0.000000 & 0.000000 & 0.000000 & 0.000000 & 1.000000\\ \midrule
        54 & 2 & 0.000000 & 0.000000 & 0.333333 & 0.666667 & 0.000000 & 57 & 2 & 0.666667 & 0.333333 & 0.000000 & 0.000000 & 0.000000 \\ \midrule
        32 & 1 & 0.000000 & 0.333333 & 0.333333 & 0.333333 & 0.000000 & 43 & 2 & 0.000000 & 1.000000 & 0.000000 & 0.000000 & 0.000000\\ \midrule
        80 & 2 & 0.100000 & 0.100000 & 0.100000 & 0.600000 & 0.100000 &  1 & 1 & 1.000000 & 0.000000 & 0.000000 & 0.000000 & 0.000000\\ \midrule
        8 & 1 & 0.333333 & 0.000000 & 0.000000 & 0.666667 & 0.000000 & 22 & 1 & 0.000000 & 0.666667 & 0.000000 & 0.000000 & 0.333333\\ \midrule
        45 & 2 & 0.000000 & 0.000000 & 0.000000 & 1.000000 & 0.000000 &   5 & 1 & 0.000000 & 0.000000 & 0.000000 & 0.000000 & 1.000000 \\ \midrule
        36 & 1 & 0.600000 & 0.100000 & 0.100000 & 0.100000 & 0.100000 & 40 & 1 & 0.100000 & 0.100000 & 0.100000 & 0.100000 & 0.600000 \\ \midrule
        24 & 1 & 0.000000 & 0.000000 & 0.666667 & 0.000000 & 0.333333 &  18 & 1 & 0.666667 & 0.000000 & 0.000000 & 0.333333 & 0.000000\\ \midrule
        27 & 1 & 0.333333 & 0.333333 & 0.000000 & 0.333333 & 0.000000 & 21 & 1 & 0.000000 & 0.666667 & 0.000000 & 0.333333 & 0.000000 \\ \midrule
        48 & 2 & 0.333333 & 0.000000 & 0.666667 & 0.000000 & 0.000000 & 58 & 2 & 0.666667 & 0.000000 & 0.333333 & 0.000000 & 0.000000 \\ 

        \bottomrule
\end{tabularx}
\end{adjustwidth}
\end{table}



\begin{table}[H]\ContinuedFloat
\renewcommand{\arraystretch}{1.3}
\tablesize{\scriptsize}
\caption{\emph{Cont.} \label{tabA1}}
\begin{adjustwidth}{-\extralength}{0cm}
		\newcolumntype{C}{>{\centering\arraybackslash}X}
		\begin{tabularx}{\fulllength}{CCCCCCCCCCCCCC}
			\noalign{\hrule height 1pt} 
			\rowcolor{green!10} 
\textbf{PN} & \textbf{R} & \textbf{Cement} & \textbf{Lime} & \textbf{Slag} & \textbf{Flyash1} & \textbf{Flyash2} & \textbf{PN} & \textbf{R} & \textbf{Cement} & \textbf{Lime} & \textbf{Slag} & \textbf{Flyash1} & \textbf{Flyash2}\\ \noalign{\hrule height .5pt} 
        56 & 2 & 0.000000 & 0.000000 & 0.000000 & 0.333333 & 0.666667 & 72 & 2 & 0.333333 & 0.000000 & 0.000000 & 0.333333 & 0.333333\\ \midrule
        64 & 2 & 0.000000 & 0.000000 & 0.666667 & 0.333333 & 0.000000 &  55 & 2 & 0.000000 & 0.000000 & 0.333333 & 0.000000 & 0.666667\\ \midrule
        26 & 1 & 0.333333 & 0.333333 & 0.333333 & 0.000000 & 0.000000 &  51 & 2 & 0.000000 & 0.333333 & 0.666667 & 0.000000 & 0.000000 \\ \midrule
        47 & 2 & 0.333333 & 0.666667 & 0.000000 & 0.000000 & 0.000000 & 20 & 1 & 0.000000 & 0.666667 & 0.333333 & 0.000000 & 0.000000\\ \midrule
        77 & 2 & 0.600000 & 0.100000 & 0.100000 & 0.100000 & 0.100000 &78 & 2 & 0.100000 & 0.600000 & 0.100000 & 0.100000 & 0.100000 \\ \midrule
        7 & 1 & 0.333333 & 0.000000 & 0.666667 & 0.000000 & 0.000000 &  4 & 1 & 0.000000 & 0.000000 & 0.000000 & 1.000000 & 0.000000 \\ \midrule
        25 & 1 & 0.000000 & 0.000000 & 0.000000 & 0.666667 & 0.333333 & 34 & 1 & 0.000000 & 0.333333 & 0.000000 & 0.333333 & 0.333333 \\ \midrule
        62 & 2 & 0.000000 & 0.666667 & 0.000000 & 0.333333 & 0.000000 & 14 & 1 & 0.000000 & 0.000000 & 0.333333 & 0.000000 & 0.666667  \\ \midrule
        12 & 1 & 0.000000 & 0.333333 & 0.000000 & 0.000000 & 0.666667 &  52 & 2 & 0.000000 & 0.333333 & 0.000000 & 0.666667 & 0.000000\\ \midrule
        60 & 2 & 0.666667 & 0.000000 & 0.000000 & 0.000000 & 0.333333 & 74 & 2 & 0.000000 & 0.333333 & 0.333333 & 0.000000 & 0.333333\\ \midrule
        23 & 1 & 0.000000 & 0.000000 & 0.666667 & 0.333333 & 0.000000  &   70 & 2 & 0.333333 & 0.000000 & 0.333333 & 0.333333 & 0.000000 \\ \midrule
        75 & 2 & 0.000000 & 0.333333 & 0.000000 & 0.333333 & 0.333333 & 10 & 1 & 0.000000 & 0.333333 & 0.666667 & 0.000000 & 0.000000\\ \midrule
        30 & 1 & 0.333333 & 0.000000 & 0.333333 & 0.000000 & 0.333333 &  59 & 2 & 0.666667 & 0.000000 & 0.000000 & 0.333333 & 0.000000 \\ 

        \bottomrule
\end{tabularx}
\end{adjustwidth}
\end{table}


\begin{adjustwidth}{-\extralength}{0cm}
 
\reftitle{{References}} 

\begin{thebibliography}{999}

\bibitem[Park \em{et~al.}(2020)Park, Lee, and Vimonsatit]{Park}
Park, H.K.; Lee, H.; Vimonsatit, V.
\newblock {Investigation of Pindan soil modified with polymer stablisers for
  road pavement}.
\newblock {\em J. Infrastruct. Preserv. Resil.} {\bf
  2020}, {\em 1},~9. https://doi.org/10.1186/s43065-020-00009-8.

\bibitem[Glossop and Wilson(1953)]{Glossop}
Glossop, R.; Wilson, G.C.
\newblock {Soil stability problems in road engineering}.
\newblock {\em Proc. Inst. Civ. Eng.} {\bf 1953},
  {\em 2},~219--253. https://doi.org/10.1680/ipeds.1953.11550.

\bibitem[Källén \em{et~al.}(2016)Källén, Heyden, Åström, and
  Lindh]{Kallenetal2016}
Källén, H.; Heyden, A.; Åström, K.; Lindh, P.
\newblock {Measuring and evaluating bitumen coverage of stones using two
  different digital image analysis methods}.
\newblock {\em Measurement} {\bf 2016}, {\em 84},~56--67. https://doi.org/10.1016/j.measurement.2016.02.007.

\bibitem[Abdullah and Al-Abdul~Wahhab(2015)]{Abdullah}
Abdullah, G.M.S.; Al-Abdul~Wahhab, H.I.
\newblock {Evaluation of foamed sulfur asphalt stabilized soils for road
  applications}.
\newblock {\em Constr. Build. Mater.} {\bf 2015}, {\em
  88},~149--158. https://doi.org/10.1016/j.conbuildmat.2015.04.013.

\bibitem[Källén \em{et~al.}(2014)Källén, Heyden, and Lindh]{Kallenetal2014}
Källén, H.; Heyden, A.; Lindh, P.
\newblock {Estimation of grain size in asphalt samples using digital image
  analysis}.
\newblock  In \emph{Proceedings of the Applications of Digital Image Processing XXXVII,} \emph{{San Diego, CA, USA, 26--29 August 2013}}; Tescher, A.G., Ed.;
  International Society for Optics and Photonics, SPIE: Bellingham, WA, USA,
  2014; Volume 9217, p. 921714. https://doi.org/10.1117/12.2061730.

\bibitem[Sudhakar \em{et~al.}(2021)Sudhakar, Duraisekaran, Vignesh, and
  Dhuvarakesh~Kanna]{Sudhakar}
Sudhakar, S.; Duraisekaran, E.; Vignesh, D.; Dhuvarakesh~Kanna, G.G.
\newblock {Performance Evaluation of Quarry Dust Treated Expansive Clay for
  Road Foundations}.
\newblock {\em Iran. J. Sci. Technol. Trans. Civ. Eng.} {\bf 2021}, {\em 45},~2637--2649. https://doi.org/10.1007/s40996-021-00645-4.

\bibitem[Nordmark \em{et~al.}(2022)Nordmark, Vestin, Hansson, and
  Kumpiene]{Nordmark}
Nordmark, D.; Vestin, J.; Hansson, L.; Kumpiene, J.
\newblock {Long-term evaluation of geotechnical and environmental properties of
  ash-stabilized road}.
\newblock {\em J. Environ. Manag.} {\bf 2022}, {\em
  318},~115504. https://doi.org/10.1016/j.jenvman.2022.115504.

\bibitem[Mitchell and {El Jack}(1966)]{Mitchell}
Mitchell, J.K.; {El Jack}, S.A.
\newblock {The Fabric of Soil---Cement and Its Formation}.
\newblock {\em Clays Clay Miner.} {\bf 1966}, {\em 14},~297--305. https://doi.org/10.1346/CCMN.1966.0140126.

\bibitem[Chen and Wang(2006)]{Chen}
Chen, H.; Wang, Q.
\newblock {The behaviour of organic matter in the process of soft soil
  stabilization using cement}.
\newblock {\em Bull. Eng. Geol. Environ.} {\bf 2006},
  {\em 65},~445--448. https://doi.org/10.1007/s10064-005-0030-1.

\bibitem[Hilt and Davidson(1960)]{Hilt}
Hilt, G.H.; Davidson, D.T.
\newblock {Lime fixation of clayey soils}.
\newblock {\em Highw. Res. Board Bull.} {\bf 1960}, {\em 262},~20--32.

\bibitem[Kaniraj and Havanagi(1999)]{Kaniraj}
Kaniraj, S.R.; Havanagi, V.G.
\newblock {Compressive strength of cement stabilized fly ash-soil mixtures}.
\newblock {\em Cem. Concr. Res.} {\bf 1999}, {\em 29},~673--677. https://doi.org/10.1016/S0008-8846(99)00018-6.

\bibitem[Markwick and Keep(1942)]{Markwick}
Markwick, A.H.D.; Keep, H.S.
\newblock {The use of low-grade aggregates and soils in the construction of
  bases for roads and aerodromes. Road Engineering Division}.
\newblock {\em  Inst. Civ. Eng. Eng. Div. Pap.}
  {\bf 1942}, {\em 1},~1--30. https://doi.org/10.1680/idivp.1942.13348.

\bibitem[Abbey \em{et~al.}(2021)Abbey, Eyo, Okeke, and Ngambi]{Abbey}
\textls[-10]{Abbey, S.J.; Eyo, E.U.; Okeke, C.A.U.; Ngambi, S.
\newblock {Experimental study on the use of RoadCem blended with by-product
  cementitious materials for stabilization of clay soils}.
\newblock {\em Constr. Build. Mater.} {\bf 2021}, {\em
  280},~122476. https://doi.org/10.1016/j.conbuildmat.2021.122476.}

\bibitem[Aldaood \em{et~al.}(2021)Aldaood, Khalil, Bouasker, and
  Al-Mukhtar]{Aldaood}
Aldaood, A.; Khalil, A.; Bouasker, M.; Al-Mukhtar, M.
\newblock {Experimental Study on the Mechanical Behavior of Cemented Soil
  Reinforced with Straw Fiber}.
\newblock {\em Geotech. Geol. Eng.} {\bf 2021}, {\em
  39},~2985--3001. https://doi.org/10.1007/s10706-020-01673-z.

\bibitem[Bell(1989)]{Bell}
Bell, F.G.
\newblock {Lime stabilization of clay soils}.
\newblock {\em Bull. Int. Assoc. Eng. Geol.} {\bf 1989}, {\em 39},~67--74. https://doi.org/10.1007/BF02592537.

\bibitem[T. \em{et~al.}(1999)T., M., and P.]{Dahlinetal1999}
{Dahlin, T.M.; Lindh,  S.P.} 
\newblock {DC Resistivity and SASW for Validation of Efficiency in Soil
  Stabilisation Prior to Road Construction}.
\newblock  In Proceedings of the 5th EEGS-ES Meeting, {Budapest, Hungary,  6--9 September} 1999; pp. 1--3. \linebreak https://doi.org/10.3997/2214-4609.201406466.

\bibitem[Nelson and Miller(1997)]{Nelson}
Nelson, J.D.; Miller, D.J.
\newblock {\em Expansive Soils: Problems and Practice in Foundation and
  Pavement Engineering}; John Wiley \& Sons Inc.: {Hoboken, NJ, USA,} 1997.

\bibitem[Jegandan \em{et~al.}(2010)Jegandan, Liska, Osman, and
  Al-Tabbaa]{Jegandanetal2010}
Jegandan, S.; Liska, M.; Osman, A.A.M.; Al-Tabbaa, A.
\newblock {Sustainable binders for soil stabilisation}.
\newblock {\em Ground Improv.} {\bf 2010}, {\em 163},~53--61. https://doi.org/10.1680/grim.2010.163.1.53.

\bibitem[Jafer(2017)]{Jafer2017}
Jafer, H.M.
\newblock Soft Soil Stabilisation Using a Novel Blended Cementitious Binder
  Produced From Waste Fly Ashes.
\newblock Ph.D. Thesis, Liverpool John Moores University, Liverpool, UK,
   2017. https://doi.org/10.24377/LJMU.t.00007541.

\bibitem[Mavroulidou \em{et~al.}(2021)Mavroulidou, Gray, Gunn, and
  Pantoja-Muñoz]{Mavroulidou}
\textls[-15]{Mavroulidou, M.; Gray, C.; Gunn, M.J.; Pantoja-Muñoz, L.
\newblock {A Study of Innovative Alkali-Activated Binders for Soil
  stabilization in the Context of Engineering Sustainability and Circular
  Economy}.
\newblock {\em Circ. Econ. Sustain.} {\textbf{2021}}. https://doi.org/10.1007/s43615-021-00112-2.}

\bibitem[Al-Swaidani \em{et~al.}(2016)Al-Swaidani, Hammoud, and
  Meziab]{AlSwaidani}
Al-Swaidani, A.; Hammoud, I.; Meziab, A.
\newblock {Effect of adding natural pozzolana on geotechnical properties of
  lime-stabilized clayey soil}.
\newblock {\em J. Rock Mech. Geotech. Eng.} {\bf
  2016}, {\em 8},~714--725. https://doi.org/10.1016/j.jrmge.2016.04.002.

\bibitem[Fasihnikoutalab \em{et~al.}(2020)Fasihnikoutalab, Pourakbar, Ball,
  Unluer, and Cristelo]{Fasihnikoutalab}
Fasihnikoutalab, M.H.; Pourakbar, S.; Ball, R.J.; Unluer, C.; Cristelo, N.
\newblock {Sustainable soil stabilization with ground granulated blast-furnace
  slag activated by olivine and sodium hydroxide}.
\newblock {\em Acta Geotech.} {\bf 2020}, {\em 15},~1981--1991. \linebreak https://doi.org/10.1007/s11440-019-00884-w.

\bibitem[Kogbara \em{et~al.}(2011)Kogbara, Yi, and Al-Tabbaa]{Kogbara}
Kogbara, R.B.; Yi, Y.; Al-Tabbaa, A.
\newblock {Process envelopes for stabilization/solidification of contaminated
  soil using lime-slag blend}.
\newblock {\em Environ. Sci. Pollut. Res.} {\bf 2011}, {\em
  18},~1286--1296. https://doi.org/10.1007/s11356-011-0480-x.

\bibitem[A.(1990)]{Verhasselt}
{Verhasselt, V.} 
\newblock {Lime-cement stabilization of wet cohesive soils}.
\newblock  In Proceedings of the 6th International Symposium on Concrete Roads, {Madrid, Spain}, {8--10 October 1990}
; pp. 67--76.

\bibitem[Sherwood(1993)]{Sherwood}
Sherwood, P.
\newblock {\em Soil Stabilization with Cement and Lime}; TRL, HMSO: London, UK,
   1993.

\bibitem[Chaddock(1996)]{Chaddock}
Chaddock, B.C.J. {The structural performance of stabilised soil in road
  foundations}.
\newblock In {\em {Lime Stabilisation}}; Thomas Telford: London, UK,  1996; pp.
  75--93.

\bibitem[Wang \em{et~al.}(2015)Wang, Wang, Jin, and Al-Tabbaa]{Wangetal2015}
Wang, F.; Wang, H.; Jin, F.; Al-Tabbaa, A.
\newblock {The performance of blended conventional and novel binders in the
  in-situ stabilization/solidification of a contaminated site soil}.
\newblock {\em J. Hazard. Mater.} {\bf 2015}, {\em 285},~46--52. https://doi.org/10.1016/j.jhazmat.2014.11.002.

\bibitem[Gandhi and Shukla(2021)]{Gandhi}
Gandhi, K.; Shukla, S.  {Enhanced Geotechnical Properties of High Plastic Clay
  Stabilized with Industrial Waste–Step for Sustainable Development.}
\newblock In {\em {Ground Improvement Techniques. Lecture Notes in Civil
  Engineering}}; Springer: Singapore,  2021; \mbox{Volume 118,} pp.~85--97. https://doi.org/10.1007/978-981-15-9988-0\_9.

\bibitem[Lindh and Lemenkova(2021)]{LindhLemenkova2021a}
Lindh, P.; Lemenkova, P.
\newblock {Evaluation of Different Binder Combinations of Cement, Slag and CKD
  for S/S Treatment of TBT Contaminated Sediments}.
\newblock {\em Acta Mech. Et Autom.} {\bf 2021}, {\em 15},~236--248. https://doi.org/10.2478/ama-2021-0030.

\bibitem[\r{A}hnberg \em{et~al.}(1995)\r{A}hnberg, Johansson, Ljungkrantz,
  Holmqvist, and Holm]{Ahnberg1995}
\r{A}hnberg, H.; Johansson, S.-E.~Retelius, A.; Ljungkrantz, C.; Holmqvist, L.;
  Holm, G.
\newblock \emph{Cement och Kalk f{\o}r Djupstabilisering av Jord (Cement and Lime For
  Deep Stabilization of Soil)};
\newblock Technical Report~48; Swedish Geotechnical Institute: Linköping,
  Sweden,  1995.\mbox{ (In Swedish)}

\bibitem[Holm and \r{A}hnberg(1987)]{HolmAhnberg}
Holm, G.; \r{A}hnberg, H.
\newblock \emph{Use of Lime-Fly Ash for Deep Soil Stabilization Purposes};
\newblock Technical Report~30; Swedish Geotechnical Institute: Linköping,
  Sweden,  1987;  pp. 59-91.
\newblock (In Swedish)

\bibitem[Bjerrum and Flodin(1960)]{Bjerrum}
Bjerrum, L.; Flodin, N.
\newblock {The Development of Soil Mechanics in Sweden, 1900-1925}.
\newblock {\em G\'{e}otechnique} {\bf 1960}, {\em 10},~1--18. https://doi.org/10.1680/geot.1960.10.1.1.

\bibitem[Sörengård \em{et~al.}(2021)Sörengård, Gago-Ferrero, B., and
  Ahrens]{Sorengard}
Sörengård, M.; Gago-Ferrero, P.; Kleja, D.B.; Ahrens, L.
\newblock {Laboratory-scale and pilot-scale stabilization and solidification
  (S/S) remediation of soil contaminated with per- and polyfluoroalkyl
  substances (PFASs)}.
\newblock {\em J. Hazard. Mater.} {\bf 2021}, {\em 402},~123453. https://doi.org/10.1016/j.jhazmat.2020.123453.

\bibitem[Lindh and Lemenkova(2022)]{LindhLemenkova2022b}
Lindh, P.; Lemenkova, P.
\newblock {Soil contamination from heavy metals and persistent organic
  pollutants (PAH, PCB and HCB) in the coastal area of Västernorrland,
  Sweden}.
\newblock {\em Gospod. Surowcami Miner. -Miner. Resour. Manag.} {\bf 2022}, {\em 38},~147--168. https://doi.org/10.24425/gsm.2022.141662.

\bibitem[{El Mahdi Safhi} \em{et~al.}(2022){El Mahdi Safhi}, Taha, {El Ghorfi},
  Hakkou, and Benzaazoua]{ElMahdi}
{El Mahdi Safhi}, A.; Taha, H.; {El Ghorfi}, M.; Hakkou, R.; Benzaazoua, M.
\newblock {Elaboration of a blended binder based on marls from phosphate mines
  waste rocks}.
\newblock {\em Constr. Build. Mater.} {\bf 2022}, {\em
  347},~128539. https://doi.org/10.1016/j.conbuildmat.2022.128539.

\bibitem[Syed \em{et~al.}(2020)Syed, GuhaRay, and Kar]{Syed}
Syed, M.; GuhaRay, A.; Kar, A.
\newblock {Stabilization of Expansive Clayey Soil with Alkali Activated
  Binders}.
\newblock {\em Geotech. Geol. Eng.} {\bf 2020}, {\em
  38},~6657--6677. https://doi.org/10.1007/s10706-020-01461-9.

\bibitem[Arulrajah \em{et~al.}(2018)Arulrajah, Yaghoubi, Disfani, Horpibulsuk,
  Bo, and Leong]{Arulrajah}
Arulrajah, A.; Yaghoubi, M.; Disfani, M.M.; Horpibulsuk, S.; Bo, M.W.; Leong,
  M.
\newblock {Evaluation of fly ash- and slag-based geopolymers for the
  improvement of a soft marine clay by deep soil mixing}.
\newblock {\em Soils Found.} {\bf 2018}, {\em 58},~1358--1370. https://doi.org/10.1016/j.sandf.2018.07.005.

\bibitem[Lindh and Lemenkova(2022)]{LindhLemenkova2022a}
Lindh, P.; Lemenkova, P.
\newblock {Geochemical tests to study the effects of cement ratio on potassium
  and TBT leaching and the pH of the marine sediments from the Kattegat Strait,
  Port of Gothenburg, Sweden}.
\newblock {\em Baltica} {\bf 2022}, {\em 35},~47--59. https://doi.org/10.5200/baltica.2022.1.4.

\bibitem[Jafer \em{et~al.}(2018)Jafer, Atherton, Sadique, Ruddock, and
  Loffill]{Jaferetal2018}
Jafer, H.M.; Atherton, W.; Sadique, M.; Ruddock, F.; Loffill, E.
\newblock {Development of a new ternary blended cementitious binder produced
  from waste materials for use in soft soil stabilization}.
\newblock {\em J. Clean. Prod.} {\bf 2018}, {\em 172},~516--528. https://doi.org/10.1016/j.jclepro.2017.10.233.

\bibitem[Eskisar(2015)]{Eskisar}
Eskisar, T.
\newblock {Influence of cement treatment on unconfined compressive strength and
  compressibility of lean clay with medium plasticity}.
\newblock {\em Arab. J. Sci. Eng.} {\bf 2015}, {\em
  40},~763--772. https://doi.org/10.1007/s13369-015-1579-z.

\bibitem[Feng \em{et~al.}(2021)Feng, Du, Zhou, Zhang, Li, Zhou, and Xia]{Feng}
Feng, Y.S.; Du, Y.J.; Zhou, A.; Zhang, M.; Li, J.S.; Zhou, S.J.; Xia, W.Y.
\newblock {Geoenvironmental properties of industrially contaminated site soil
  solidified/stabilized with a sustainable by-product-based binder}.
\newblock {\em Sci. Total Environ.} {\bf 2021}, {\em 765},~142778. https://doi.org/10.1016/j.scitotenv.2020.142778.

\bibitem[Rodriguez \em{et~al.}(1988)Rodriguez, {del Castillo}, and
  Sowers]{Rodriguez}
Rodriguez, R.; {del Castillo}, H.; Sowers, G.F.
\newblock {\em Soil Mechanics in Highway Engineering}, 2nd ed.; Series on Rock
  and Soil Mechanics, Trans Tech Publications;  {Elsevier:  Amsterdam, The Netherlands, }1988.
 
\bibitem[Christensen(1969)]{Christensen}
Christensen, A.P. {Cement Modification of Clay Soils}.
\newblock In {\em {RD002}}; Portland Cement Association: Skokie, IL, USA,
  1969; pp. 1--14.

\bibitem[Alshkane \em{et~al.}(2020)Alshkane, Rashed, and Daoud]{Alshkane}
Alshkane, Y.M.; Rashed, K.A.; Daoud, H.S.
\newblock {Unconfined Compressive Strength (UCS) and Compressibility Indices
  Predictions from Dynamic Cone Penetrometer Index (DCP) for Cohesive Soil in
  Kurdistan Region/Iraq}.
\newblock {\em Geotech. Geol. Eng.} {\bf 2020}, {\em
  38},~3683--3695. https://doi.org/10.1007/s10706-020-01245-1.

\bibitem[Onyelowe \em{et~al.}(2022)Onyelowe, Jalal, Iqbal, Rehman, and
  Ibe]{Onyelowe}
Onyelowe, K.C.; Jalal, F.E.; Iqbal, M.; Rehman, Z.U.; Ibe, K.
\newblock {Intelligent modeling of unconfined compressive strength (UCS) of
  hybrid cement-modified unsaturated soil with nanostructured quarry fines
  inclusion}.
\newblock {\em Innov. Infrastruct. Solut.} {\bf 2022}, {\em 7},~98. https://doi.org/10.1007/s41062-021-00682-y.

\bibitem[Tabarsa \em{et~al.}(2021)Tabarsa, Latifi, Osouli, and
  Bagheri]{Tabarsa}
Tabarsa, A.; Latifi, N.; Osouli, A.; Bagheri, Y.
\newblock {Unconfined compressive strength prediction of soils stabilized using
  artificial neural networks and support vector machines}.
\newblock {\em Front. Struct. Civ. Eng.} {\bf 2021}, {\em
  15},~520--536. https://doi.org/10.1007/s11709-021-0689-9.

\bibitem[Sihag \em{et~al.}(2021)Sihag, Suthar, and Mohanty]{Sihag}
Sihag, P.; Suthar, M.; Mohanty, S.
\newblock {Estimation of UCS-FT of Dispersive Soil Stabilized with Fly Ash,
  Cement Clinker and GGBS by Artificial Intelligence}.
\newblock {\em Iran. J. Sci. Technol. Trans. Civ. Eng.} {\bf 2021}, {\em 45},~901--912. https://doi.org/10.1007/s40996-019-00329-0.

\bibitem[Liu \em{et~al.}(2015)Liu, Shao, Xu, and Wu]{Liuetal2015}
Liu, Z.; Shao, J.; Xu, W.; Wu, Q.
\newblock {Indirect estimation of unconfined compressive strength of carbonate
  rocks using extreme learning machine}.
\newblock {\em Acta Geotech.} {\bf 2015}, {\em 10},~651--663. https://doi.org/10.1007/s11440-014-0316-1.

\bibitem[Furtney \em{et~al.}(2022)Furtney, Thielsen, Fu, and Le~Goc]{Furtney}
Furtney, J.K.; Thielsen, C.; Fu, W.; Le~Goc, R.
\newblock {Surrogate Models in Rock and Soil Mechanics: Integrating Numerical
  Modeling and Machine Learning}.
\newblock {\em Rock Mech. Rock Eng.} {\bf 2022}, {\em
  55},~2845--2859. https://doi.org/10.1007/s00603-021-02720-8.

\bibitem[Gheibi and Hedayat(2018)]{Gheibi}
Gheibi, A.; Hedayat, A.
\newblock {Ultrasonic investigation of granular materials subjected to
  compression and crushing}.
\newblock {\em Ultrasonics} {\bf 2018}, {\em 87},~112--125. https://doi.org/10.1016/j.ultras.2018.02.006.

\bibitem[Jiang \em{et~al.}(2021)Jiang, Zhan, Zhang, Jiang, Zhuang, and
  Fan]{Jiangetal2021}
Jiang, H.; Zhan, H.; Zhang, J.; Jiang, R.; Zhuang, C.; Fan, P.
\newblock {Detecting Stress Changes and Damage in Full-Size Concrete T-Beam and
  Slab with Ultrasonic Coda Waves}.
\newblock {\em J. Struct. Eng.} {\bf 2021}, {\em
  147},~04021140. https://doi.org/10.1061/(ASCE)ST.1943-541X.0003090.

\bibitem[Jiang \em{et~al.}(2017)Jiang, Zhang, and Jiang]{Jiangetal2017}
Jiang, H.; Zhang, J.; Jiang, R.
\newblock {Stress Evaluation for Rocks and Structural Concrete Members through
  Ultrasonic Wave Analysis: Review}.
\newblock {\em J. Mater. Civ. Eng.} {\bf 2017}, {\em
  29},~04017172. https://doi.org/10.1061/(ASCE)MT.1943-5533.0001935.

\bibitem[Lindh and Lemenkova(2021)]{LindhLemenkova2021b}
Lindh, P.; Lemenkova, P.
\newblock {Resonant Frequency Ultrasonic P-Waves for Evaluating Uniaxial
  Compressive Strength of the Stabilized Slag–Cement Sediments}.
\newblock {\em Nord. Concr. Res.} {\bf 2021}, {\em 65},~39--62. https://doi.org/10.2478/ncr-2021-0012.

\bibitem[Lacidogna \em{et~al.}(2019)Lacidogna, Piana, and
  Carpinteri]{Lacidogna}
Lacidogna, G.; Piana, G.; Carpinteri, A.
\newblock {Damage monitoring of three-point bending concrete specimens by
  acoustic emission and resonant frequency analysis}.
\newblock {\em Eng. Fract. Mech.} {\bf 2019}, {\em 210},~203--211. https://doi.org/10.1016/j.engfracmech.2018.06.034.

\bibitem[Lindh and Lemenkova(2022)]{LindhLemenkova2022c}
Lindh, P.; Lemenkova, P.
\newblock {Seismic velocity of P-waves to evaluate strength of stabilized soil
  for Svenska Cellulosa Aktiebolaget Biorefinery Östrand AB, Timrå}.
\newblock {\em Bull. Pol. Acad. Sci. Tech. Sci.}
  {\bf 2022}, {\em 70},~e141593. https://doi.org/10.24425/bpasts.2022.141593.

\bibitem[Lin and Wang(2020)]{LinWang2020}
Lin, S.; Wang, Y.
\newblock {Crack-Depth Estimation in Concrete Elements Using Ultrasonic
  Shear-Horizontal Waves}.
\newblock {\em J. Perform. Constr. Facil.} {\bf 2020},
  {\em 34},~04020064. https://doi.org/10.1061/(ASCE)CF.1943-5509.0001473.

\bibitem[Kurtulus \em{et~al.}(2010)Kurtulus, Sertcelik, Canbay, and
  Sertcelik]{Kurtulus}
Kurtulus, C.; Sertcelik, F.; Canbay, M.M.; Sertcelik, I.
\newblock {Estimation of Atterberg limits and bulk mass density of an expansive
  soil from P-wave velocity measurements}.
\newblock {\em Bull. Eng. Geol. Environ.} {\bf 2010},
  {\em 69},~153--154. https://doi.org/10.1007/s10064-009-0237-7.

\bibitem[Ryden \em{et~al.}(2006)Ryden, Ekdahl, and Lindh]{Ryden}
Ryden, N.; Ekdahl, U.; Lindh, P.
\newblock {Quality Control of Cement stabilized Soil Using Non-destructive
  Seismic Tests}.
\newblock  In Proceedings of the Advanced Testing of Fresh Cementitious Materials, {Stuttgart, Germany, 3--4 August} 2006; {%MDPI: Please confirm if this information can be removed. -- Answer: ok, agree, we removed it.
}; pp. 1--5. 

\bibitem[Sarro \em{et~al.}(2021)Sarro, Assis, and Ferreira]{Sarro}
Sarro, W.S.; Assis, G.M.; Ferreira, G.C.S.
\newblock {Experimental investigation of the UPV wavelength in compacted soil}.
\newblock {\em Constr. Build. Mater.} {\bf 2021}, {\em
  272},~121834. https://doi.org/10.1016/j.conbuildmat.2020.121834.

\bibitem[Shi \em{et~al.}(2017)Shi, Wen, and Meng]{Shietal2017}
Shi, Z.; Wen, Y.; Meng, Q.
\newblock {Propagation Attenuation of Plane Waves in Saturated Soil by Pile
  Barriers}.
\newblock {\em Int. J. Geomech.} {\bf 2017}, {\em
  17},~04017053. https://doi.org/10.1061/(ASCE)GM.1943-5622.0000963.

\bibitem[{SSAB Merox AB}(2008)]{SSAB}
{SSAB Merox AB}.
\newblock {Merit 5000 Produktspecifikation, 2012 Doseringsråd---Användning
  Av Merit 5000 I Betong Tillämpning Av SS-EN 206-1 OCH SS 13 70 03:2008
  (Merit 5000 Product Specification, 2012 Dosage Advice---Use of Merit 5000 in
  Concrete Application)}. 2008.
\newblock Available online:  \url{https://docplayer.se/8579935-Doseringsrad-merit-5000.html}  (accessed on 3 November 2022 %MDPI: Please add the access date (format: Date Month Year), e.g., accessed on 1 January 2020. The following highlights are the same. -- Answer: ok, corrected. We checked all the web links today, so we have inserted "accessed on 3 November 2022" where required.
).

\bibitem[SIS(2003)]{SIS2003}
SIS.
\newblock {Obundna Och Hydrauliskt Bundna VäGmaterial---Del 41:
  Provningsmetod FöR BestäMning Av TryckhåLlfastheten FöR Hydrauliskt
  Bundna Material (Unbound and Hydraulically Bound Mixtures---Part 41: Test
  Method for the Determination of the Compressive Strength of Hydraulically
  Bound Mixtures)}; SS-EN 13286-41. 2003. 
\newblock Available online:  \url{https://www.sis.se/api/document/preview/33822/} (accessed on 3 November 2022).

\bibitem[SIS(2017)]{SIS2018}
\emph{ISO 17892-7:2017}; Geotechnical Investigation and Testing---Laboratory Testing of Soil---Part 7: Unconfined Compression Test. {SIS:  Stockholm, Sweden}, 2017. 
\newblock Available online: \url{https://www.iso.org/standard/70789.html} (accessed on 3 November 2022).

\bibitem[SIS(2014)]{SIS2014}
\emph{ISO 17892-1:2014}; Geotechnical Investigation and Testing---Laboratory Testing of Soil---Part 1: Determination of Water Content. {SIS:  Stockholm, Sweden}, 2014. Available online: \url{https://www.sis.se/api/document/preview/104733/} (accessed on 3 November 2022).
 

\bibitem[SIS(2005)]{SIS2005}
\emph{SS-EN 137244}; Betongprovning---HåRdnad Betong---Avflagning Vid Frysning
  (Concrete Testing---Hardened Concrete---Flaking During Freezing.  {SIS:  Stockholm, Sweden}, 2005.  Available online: 
  \url{https://www.sis.se/produkter/byggnadsmaterial-och-byggnader/byggnadsmaterial/betong-och-betongprodukter/ss1372442005/}  (accessed on 3 November 2022).

\bibitem[Myers and Montgomery(1995)]{MyersMontgomery}
Myers, R.H.; Montgomery, D.C.
\newblock {\em Response Surface Methodology}; Wiley: {Hoboken, NJ, USA}, 1995.

\bibitem[Montgomery(1996)]{Montgomery}
Montgomery, D.C.
\newblock {\em Design and Analysis of Experiments}, 7th ed.; Wiley: {Hoboken, NJ, USA}, 1996.

\bibitem[Lindh(2001)]{Lindh2001}
Lindh, P.
\newblock {Optimizing binder blends for shallow stabilisation of fine-grained
  soils}.
\newblock {\em Ground Improv.} {\bf 2001}, {\em 5},~23–34. https://doi.org/10.1680/grim.2001.5.1.23.

\bibitem[Lindh(2004)]{Lindh2004}
Lindh, P.
\newblock Compaction---And Strength Properties of Stabilised and Unstabilised
  Fine-Grained Tills.
\newblock Ph.D. Thesis, Lund University, Lund, Sweden,  2004. https://doi.org/10.13140/RG.2.1.1313.6481.

\bibitem[Smith(2005)]{Smith}
Smith, W.F.
\newblock {\em Experimental Design for Formulation}; SIAM: {Philadelphia, PA, USA}; ASA: {Washington, DC, USA,} 2005.

\end{thebibliography}

\section*{Short Biography of Authors}
\bio
{\raisebox{-0.35cm}{\includegraphics[width=3.5cm,height=5.3cm,clip,keepaspectratio]{Photo_author_1.jpg}}}
{\textbf{Polina Lemenkova} is an Engineer and a PhD student in the Image Analysis Research Unit (LISA-IA) in École polytechnique de Bruxelles, Université Libre de Bruxelles. Her research interests are focused on image processing, cartography, spatial data analysis, remote sensing and applied programming. She has a strong research experience in Earth and environmental sciences, mapping and data visualization, specifically in geological and geophysical domains. Her current topic is on seismic data processing using machine learning methods and scripting languages.}
%
\bio
{\raisebox{-0.35cm}{\includegraphics[width=3.5cm,height=5.3cm,clip,keepaspectratio]{Photo_author_2.jpg}}}
{\textbf{Per Lindh} is a Senior Specialist in Swedish Transport Administration and an Adjunct Senior Lecturer at Lund University. He is a Civil Engineer with a sound experience in soil stabilization with cementitious binders and stabilization of contaminated soils and sediments using statistical methods. His area of expertise is focused on soil mechanics, compressive strength, measurements of materials, sustainable ground constructions and building, including nondestructive testing of material properties by seismic and sonic methods. He actively participates in geotechnical projects on highway and pavement engineering and solidification of soil for asphalt and road constructions.}

%%%%%%%%%%%%%%%%%%%%%%%%%%%%%%%%%%%%%%%%%%

\end{adjustwidth}

%=====================================
% References: external bibliography
%=====================================
%\externalbibliography{yes}
%\bibliography{DynamicsStrength.bib}
%\nocite{*}

\end{document}

